\documentclass[11pt]{article}
\usepackage[T1]{fontenc}
\usepackage[utf8]{inputenc}
\usepackage{lmodern}
\usepackage{microtype}
\usepackage[a4paper,margin=30mm,headheight=15pt]{geometry}
\usepackage{amsmath,amssymb,amsthm,mathtools}
\usepackage{mathrsfs}
\usepackage{bm}
\usepackage{booktabs,longtable,array}
\usepackage{enumitem}
\usepackage{xcolor}
\usepackage{fancyhdr}
\usepackage{hyperref}
\usepackage[nameinlink,capitalise,noabbrev]{cleveref}
\usepackage{tikz}
\usetikzlibrary{arrows.meta,positioning,calc}

\definecolor{linkblue}{RGB}{25,73,133}
\definecolor{darkred}{RGB}{145,35,35}
\hypersetup{
  colorlinks=true,
  linkcolor=linkblue,
  citecolor=darkred,
  urlcolor=linkblue,
  pdftitle={The Fabius Function and Rvachev's up-Function},
  pdfauthor={Human-readable exposition based on the Lean development in ProveIt}
}

\pagestyle{fancy}
\fancyhf{}
\fancyhead[L]{The Fabius function and Rvachev's $\operatorname{up}$-function}
\fancyhead[R]{\thepage}
\renewcommand{\headrulewidth}{0.4pt}

\numberwithin{equation}{section}
\setcounter{tocdepth}{2}
\setcounter{secnumdepth}{3}
\setlist{itemsep=2pt,topsep=5pt}
\allowdisplaybreaks

\newtheorem{theorem}{Theorem}[section]
\newtheorem{lemma}[theorem]{Lemma}
\newtheorem{proposition}[theorem]{Proposition}
\newtheorem{corollary}[theorem]{Corollary}
\newtheorem{conjecture}[theorem]{Conjecture}
\theoremstyle{definition}
\newtheorem{definition}[theorem]{Definition}
\newtheorem{example}[theorem]{Example}
\newtheorem{algorithm}[theorem]{Algorithm}
\theoremstyle{remark}
\newtheorem{remark}[theorem]{Remark}
\newtheorem{warning}[theorem]{Warning}

\DeclareMathOperator{\supp}{supp}
\DeclareMathOperator{\den}{den}
\DeclareMathOperator{\wt}{wt}
\DeclareMathOperator{\Res}{Res}
\DeclareMathOperator{\Var}{Var}
\DeclareMathOperator{\Log}{Log}
\DeclareMathOperator{\sgn}{sgn}
\newcommand{\R}{\mathbb{R}}
\newcommand{\C}{\mathbb{C}}
\newcommand{\Q}{\mathbb{Q}}
\newcommand{\Z}{\mathbb{Z}}
\newcommand{\N}{\mathbb{N}}
\newcommand{\dd}{\,\mathrm{d}}
\newcommand{\e}{\mathrm{e}}
\newcommand{\ii}{\mathrm{i}}
\newcommand{\upf}{\operatorname{up}}
\newcommand{\sinc}{\operatorname{sinc}}
\newcommand{\eps}{\varepsilon}
\newcommand{\Th}{\Theta}
\newcommand{\half}{\tfrac12}
\newcommand{\Ltwo}{\log 2}
\newcommand{\binomtri}[1]{\binom{#1}{2}}
\newcommand{\qbinom}[2]{\genfrac{[}{]}{0pt}{}{#1}{#2}_{1/2}}

\title{\bfseries The Fabius Function and Rvachev's \texorpdfstring{$\operatorname{up}$}{up}-Function\\[4pt]
\large Analysis, exact dyadic arithmetic, and the corrected full endpoint asymptotic expansion}
\author{A human-readable exposition based on the Lean formalization\\
\normalsize in Vladimir Reshetnikov's \texttt{ProveIt/Analysis/FabiusFunction}}
\date{24 August 2026}

\begin{document}
\maketitle

\begin{abstract}
The Fabius function is the unique smooth distribution function on the unit
interval satisfying the dyadic differential equation $F'(x)=2F(2x)$ on the
left half of the interval and the reflection law $F(1-x)=1-F(x)$.  Its even,
compactly supported twin is Rvachev's $\upf$-function, the distinguished
$C^\infty$ bump on $[-1,1]$ satisfying
\[
 \upf'(x)=2\bigl(\upf(2x+1)-\upf(2x-1)\bigr),
 \qquad \upf(0)=1.
\]
A third, signed global extension satisfies the simpler scaling equation
$\Th'(x)=2\Th(2x)$.  Keeping these three objects separate resolves a number
of notational ambiguities in the literature.

This article develops the principal results formalized in Lean in the
\texttt{ProveIt} repository.  We give constructions by infinite convolution,
random series, Fourier products, and step functions; prove smoothness,
positivity, support, partition-of-unity, Poisson, moment, derivative, and
non-analyticity results; derive exact rational formulas at every dyadic point;
explain the fast bit-recursive evaluator and the finite
$q$-binomial--Thue--Morse formula; and treat the associated integer,
denominator, and parity phenomena.

Particular attention is paid to the endpoint $x\downarrow0$.  The logarithm
of $F(x)$ has a quadratic logarithmic scale, but its sharp main term contains
a genuine one-periodic fluctuation.  We derive its exact mean and its
absolutely convergent Fourier series with Gamma--zeta coefficients, introduce
the exact lower-Lambert phase $\lambda2^{-\lambda}=x$, and state the full
all-orders expansion in inverse powers of $\lambda$.  The first correction is
given explicitly.  This also explains why a commonly repeated purely
elementary asymptotic, with the periodic term omitted, cannot have a vanishing
error.
\end{abstract}

\tableofcontents
\newpage

\section{Introduction and conventions}

The Fabius function is remarkable for combining features that normally live
in rather different parts of analysis.  It is $C^\infty$ but nowhere
real-analytic on its active interval; it is governed by a self-similar
functional-differential equation; its values at dyadic rationals are exactly
rational; its Fourier transform is an elementary infinite product; its
moments conceal striking integer sequences; and its endpoint behavior is
controlled by a Lambert--$W$ saddle together with a tiny but indispensable
periodic oscillation.

The name ``Fabius function'' is used in two incompatible ways.  Some authors
mean a distribution function $F$ on $[0,1]$, while others use the same letter
for an unbounded signed solution of $f'(x)=2f(2x)$.  The compactly supported
Rvachev function is then introduced by a shift and a reflection, sometimes
without a change of notation.  We shall instead use three symbols throughout:
\begin{center}
\begin{tabular}{@{}lll@{}}
\toprule
symbol & object & normalization/domain convention\\
\midrule
$F$ & bounded Fabius function & $F=0$ on $(-\infty,0]$, $F=1$ on $[1,\infty)$\\
$\upf$ & Rvachev bump & even, supported on $[-1,1]$, $\upf(0)=1$\\
$\Th$ & signed global extension & $\Th=F$ on $[0,1]$, $\Th'=2\Th(2\cdot)$\\
\bottomrule
\end{tabular}
\end{center}
This convention follows the formal development and prevents identities valid
for $\Th$ from being mistakenly asserted for the bounded function $F$.

We use the Fourier transform
\begin{equation}\label{eq:fourier-convention}
 \widehat g(z)=\int_{\R}g(t)\e^{-2\pi\ii zt}\dd t.
\end{equation}
The binary weight $\wt(n)$ is the number of $1$'s in the binary expansion of
$n$, and
\begin{equation}\label{eq:tm-sign}
 \epsilon_n=(-1)^{\wt(n)}
\end{equation}
is the Thue--Morse sign.  We write $L=\log2$.  The Stieltjes constants are
normalized by
\[
 \zeta(1+z)=\frac1z+\sum_{m\ge0}\frac{(-1)^m\gamma_m}{m!}z^m,
 \qquad \gamma_0=\gamma.
\]

\subsection{A guide to the main results}

The structural equivalence between $F$ and $\upf$ is established in
\Cref{sec:three-functions}.  The random-series and Fourier constructions are
in \Cref{sec:construction}, and the basic analytic properties are in
\Cref{sec:basic-properties}.  The global extension and its derivative tower
appear in \Cref{sec:global-extension}.  Moments and generating functions are
handled in \Cref{sec:moments}.

The exact arithmetic begins in \Cref{sec:dyadic-closed-form}.  The central
formula is \Cref{thm:closed-dyadic}, which evaluates every dyadic value by a
finite Thue--Morse sum and finitely many even moments.  A second evaluator,
based on terminating Taylor recursion and the binary digits of the numerator,
is given in \Cref{sec:bit-recursion}.  The much less obvious
$q$-binomial formula is proved in \Cref{sec:qbinomial}.  Integrality,
denominator bounds, and parity are collected in \Cref{sec:arithmetic}.

The endpoint analysis begins in \Cref{sec:asymptotics}.  The exact dyadic
renormalization of the negative Laplace product produces a periodic function
$\Psi$ with explicit Gamma--zeta Fourier coefficients.  The lower-Lambert
phase is introduced in \Cref{sec:lambert-phase}; the sharp main theorem is
\Cref{thm:sharp-main}; and the all-orders expansion is
\Cref{thm:full-expansion}.  The corrected elementary and Lambert--$W$ forms
are displayed in \Cref{sec:corrected-explicit}.

\section{The three related functions}\label{sec:three-functions}

\subsection{The bounded Fabius function}

\begin{definition}\label{def:bounded-fabius}
A \emph{bounded Fabius function} is a function $F:\R\to[0,1]$ such that
\begin{enumerate}[label=\textup{(\roman*)}]
\item $F(x)=0$ for $x\le0$ and $F(x)=1$ for $x\ge1$;
\item $F\in C^\infty(\R)$;
\item $F(1-x)=1-F(x)$ for $0\le x\le1$;
\item $F'(x)=2F(2x)$ for $0\le x\le\half$.
\end{enumerate}
\end{definition}

The endpoint constancy and smoothness force every derivative of positive
order to vanish at $0$ and $1$.  The reflection law gives $F(\half)=\half$.
Differentiating the reflection identity and using the left-half equation
shows that on the right half
\begin{equation}\label{eq:F-right-derivative}
 F'(x)=2F(2-2x),\qquad \half\le x\le1.
\end{equation}
Thus $F$ is increasing, with a symmetric density about $1/2$.

\subsection{Rvachev's compactly supported twin}

\begin{definition}\label{def:up-from-F}
Given $F$, define
\begin{equation}\label{eq:up-from-F}
 \upf(x)=
 \begin{cases}
  F(x+1),&x\le0,\\
  F(1-x),&x\ge0.
 \end{cases}
\end{equation}
\end{definition}

The two pieces agree to all orders at $0$.  Indeed they have the common value
$1$, and their positive-order derivatives are related by the symmetry of the
density.  The function is even, supported on $[-1,1]$, positive in the open
interval, and normalized by $\upf(0)=1$.

\begin{theorem}[Rvachev equation]\label{thm:rvachev-equation}
The function $\upf$ satisfies
\begin{equation}\label{eq:up-diffeq}
 \upf'(x)=2\bigl(\upf(2x+1)-\upf(2x-1)\bigr)
 \qquad(x\in\R).
\end{equation}
Conversely, there is a unique even $C^\infty$ function supported on $[-1,1]$
with integral $1$ that satisfies \eqref{eq:up-diffeq}; it has $\upf(0)=1$ and
gives the bounded Fabius function by
\begin{equation}\label{eq:F-from-up}
 F(x)=\upf(x-1),\qquad 0\le x\le1,
\end{equation}
with constant continuation outside the unit interval.
\end{theorem}

\begin{proof}
For $-1<x<0$, \eqref{eq:up-from-F} and the two half-interval formulas for
$F'$ give
\[
 \upf'(x)=F'(x+1)=2\upf(2x+1),
\]
because $2x-1<-1$ and hence $\upf(2x-1)=0$.  The case $0<x<1$ follows by
evenness, and outside $[-1,1]$ both sides vanish.  Smoothness extends the
identity through the finitely many boundary points.

Existence and uniqueness are proved below by Fourier transform.  Once
$\int\upf=1$ is known, integration of \eqref{eq:up-diffeq} from $-1$ to $0$
gives
\[
 \upf(0)=2\int_{-1}^{0}\upf(2t+1)\dd t
       =\int_{-1}^{1}\upf(u)\dd u=1.
\]
On $[-1,0]$ the same integration gives
\begin{equation}\label{eq:up-left-integral}
 \upf(x)=2\int_{-1}^{x}\upf(2t+1)\dd t.
\end{equation}
It follows that $F(x)=\upf(x-1)$ obeys $F'(x)=2\upf(2x-1)=2F(2x)$ on the
left half.  The partition identity proved in \Cref{prop:partition-unity}
gives the reflection law.  Hence the two descriptions are equivalent.
\end{proof}

\subsection{The signed global extension}

The simplest differential equation belongs not to the bounded $F$ but to a
signed global function.  Define
\begin{equation}\label{eq:theta-def-intro}
 \Th(x)=\sum_{m=0}^{\infty}\epsilon_m\upf(x-2m-1).
\end{equation}
Because $\upf$ has compact support, the sum is locally finite.  On the first
unit interval only the term $m=0$ is present, and therefore
\begin{equation}\label{eq:theta-agrees-F}
 \Th(x)=\upf(x-1)=F(x),\qquad 0\le x\le1.
\end{equation}
The function $\Th$ is zero on the negative half-line, is alternately reflected
and signed on consecutive blocks of length $2$, and is not bounded by the
right-hand constant convention that defines $F$.

\begin{warning}
The equation $f'(x)=2f(2x)$ is valid globally for $\Th$, not for the bounded
function $F$.  On $[0,\half]$ they agree and hence the familiar Fabius
equation holds there, but for $x\ge1$ the bounded function is constant while
$\Th$ continues to oscillate in signed compact arcs.
\end{warning}

\section{Construction by probability, convolution, and Fourier transform}
\label{sec:construction}

\subsection{An infinite random series}

Let $V_1,V_2,\ldots$ be independent random variables uniformly distributed
on $[-1,1]$, and set
\begin{equation}\label{eq:Y-random-series}
 Y=\sum_{j=1}^{\infty}\frac{V_j}{2^j}.
\end{equation}
The series converges absolutely almost surely and $|Y|\le1$.  Its
characteristic function, in the convention \eqref{eq:fourier-convention}, is
\begin{equation}\label{eq:char-product}
 \mathbb E\e^{-2\pi\ii zY}
 =\prod_{j=1}^{\infty}
   \frac{\sin(2\pi z/2^j)}{2\pi z/2^j}
 =\prod_{h=0}^{\infty}\frac{\sin(\pi z/2^h)}{\pi z/2^h}.
\end{equation}
The product will be denoted by $\Phi(z)$.

\begin{lemma}[Rapid decay of the product]\label{lem:rapid-product}
For every integer $N\ge0$ there is $C_N$ such that
\[
 |\Phi(\xi)|\le C_N(1+|\xi|)^{-N},\qquad \xi\in\R.
\]
The same is true after multiplying by any fixed power of $\xi$.
\end{lemma}

\begin{proof}
The elementary estimate $|\sin u/u|\le\min(1,|u|^{-1})$ implies that, once
$|\xi|$ is large, any chosen $N$ initial factors for which
$|\pi\xi/2^h|\ge1$ contribute at most a constant times $|\xi|^{-N}$; the
remaining factors have modulus at most $1$.  The range of small $\xi$ is
absorbed into the constant.  Taking more factors proves the second statement.
\end{proof}

Fourier inversion therefore gives a Schwartz-class density
\begin{equation}\label{eq:up-fourier-inversion}
 \upf(x)=\int_{\R}\Phi(\xi)\e^{2\pi\ii x\xi}\dd\xi.
\end{equation}
The law of $Y$ is supported on $[-1,1]$, so the density vanishes outside this
interval.  Symmetry of $Y$ makes $\upf$ even, and $\int\upf=\Phi(0)=1$.
Since $\Phi$ is rapidly decreasing, differentiation under the integral shows
that $\upf\in C^\infty$.

\begin{theorem}[Existence and uniqueness]\label{thm:existence-uniqueness}
The density in \eqref{eq:up-fourier-inversion} is the unique normalized
$C^\infty$ solution of the Rvachev equation.  Consequently the bounded Fabius
function exists and is unique.
\end{theorem}

\begin{proof}
The product satisfies
\begin{equation}\label{eq:fourier-refinement}
 \Phi(z)=\frac{\sin\pi z}{\pi z}\Phi(z/2).
\end{equation}
Taking inverse Fourier transforms yields \eqref{eq:up-diffeq}; alternatively,
one can obtain the equation directly by conditioning on $V_1$ and
differentiating the resulting convolution integral.

Conversely, if $g$ is an integrable normalized solution of
\eqref{eq:up-diffeq}, Fourier transformation gives
\[
 \widehat g(z)=\frac{\sin\pi z}{\pi z}\widehat g(z/2).
\]
Iterating $n$ times gives
\[
 \widehat g(z)=\prod_{h=0}^{n-1}\frac{\sin(\pi z/2^h)}{\pi z/2^h}
                \widehat g(z/2^n).
\]
Since $\widehat g(0)=1$ and $\widehat g$ is continuous at the origin, letting
$n\to\infty$ gives $\widehat g=\Phi$.  Injectivity of the Fourier transform
then gives $g=\upf$.
\end{proof}

\subsection{Infinite convolution}

Let $b_a(x)=(2a)^{-1}\mathbf1_{[-a,a]}(x)$ be the centered uniform density.
Then
\begin{equation}\label{eq:infinite-convolution}
 \upf=b_{1/2}*b_{1/4}*b_{1/8}*\cdots
\end{equation}
in the sense of weak convergence of probability measures, and in fact in
$C^m$ after enough factors for each fixed $m$.  Formula
\eqref{eq:infinite-convolution} makes nonnegativity immediate.  It also gives
an efficient conceptual proof of compact support: the Minkowski sum of the
supports has radius $\sum_{j\ge1}2^{-j}=1$.

Strict positivity in $(-1,1)$ can be seen probabilistically.  Given
$|x|<1$, choose finitely many coordinates of the random series in small open
intervals whose weighted sum lies near $x$; the remaining tail has arbitrarily
small support.  The resulting cylinder event has positive probability, so
every sufficiently small interval about $x$ has positive mass.  Continuity
and the differential identity then rule out an interior zero.  Equivalently,
\eqref{eq:up-left-integral}, evenness, and an elementary propagation argument
show that positivity at one point propagates to the whole open support.

\subsection{The distribution-function interpretation}

Let $U_1,U_2,\ldots$ be independent uniform random variables on $[0,1]$ and
put
\begin{equation}\label{eq:X-random-series}
 X=\sum_{j=1}^{\infty}\frac{U_j}{2^j}.
\end{equation}
Writing $V_j=2U_j-1$ gives $Y=2X-1$.  The following identity is one of the
most useful ways to understand the Fabius function.

\begin{theorem}[Probability representation]\label{thm:probability-F}
For $0\le x\le1$,
\begin{equation}\label{eq:F-CDF}
 F(x)=\mathbb P\{X\le x\}=\upf(x-1).
\end{equation}
The density of $X$ is
\begin{equation}\label{eq:F-density}
 F'(x)=2\upf(2x-1).
\end{equation}
\end{theorem}

\begin{proof}
Since $Y=2X-1$, the density transformation gives the density
$2\upf(2x-1)$ for $X$.  Hence its distribution function is
\[
 \int_0^x2\upf(2t-1)\dd t=\int_{-1}^{2x-1}\upf(u)\dd u.
\]
On the other hand, integrating the Rvachev equation on $[-1,x-1]$ gives
\[
 \upf(x-1)=2\int_{-1}^{x-1}\upf(2t+1)\dd t
          =\int_{-1}^{2x-1}\upf(u)\dd u.
\]
Thus the two expressions agree.
\end{proof}

The support and symmetry of $X$ are now transparent.  Since replacing every
$U_j$ by $1-U_j$ replaces $X$ by $1-X$, the distribution is symmetric about
$1/2$, which immediately gives $F(1-x)=1-F(x)$.

\section{Step approximants and elementary structure}\label{sec:basic-properties}

\subsection{A sequence of combinatorial step functions}

Define polynomials $p_n$ by
\begin{equation}\label{eq:p-recurrence}
 p_0(X)=1,\qquad p_n(X)=p_{n-1}(X^2)(1+X)^n.
\end{equation}
A simple induction gives the two useful product forms
\begin{equation}\label{eq:p-products}
 p_n(X)=\prod_{j=1}^{n}\frac{1-X^{2^j}}{1-X}
       =\prod_{j=1}^{n}(1+X+\cdots+X^{2^j-1}).
\end{equation}
Write
\[
 p_n(X)=\sum_{m=0}^{g_n}a_{n,m}X^m.
\]
Then $g_0=0$, $g_n=2g_{n-1}+n$, and
\begin{equation}\label{eq:degree-limit}
 \frac{g_n}{2^n}=\sum_{j=1}^{n}\frac{j}{2^j}\longrightarrow2.
\end{equation}
The coefficient $a_{n,m}$ counts representations
\[
 m=s_1+\cdots+s_n,
 \qquad 0\le s_j\le2^j-1.
\]

For an interval $I=[a,b]$, let $\mathbf1_I^*$ be $1$ in its interior,
$1/2$ at either endpoint, and $0$ elsewhere.  The half-weight convention is
important: adjacent intervals share endpoints, and using ordinary closed
indicators would double-count those points.  Put
\begin{equation}\label{eq:step-approx}
 \upf_n(x)=2^{n-\binom{n+1}{2}}
 \sum_{m=0}^{g_n}a_{n,m}\,
 \mathbf1^*_{I_{n,m}}(x),
\end{equation}
where
\begin{equation}\label{eq:step-interval}
 I_{n,m}=\left[
 \frac{2m-1-g_n}{2^{n+1}},
 \frac{2m+1-g_n}{2^{n+1}}
 \right].
\end{equation}

\begin{theorem}[Step-function approximation]\label{thm:step-approx}
The functions $\upf_n$ are even, unimodal, satisfy $\upf_n(0)=1$, and
converge pointwise to $\upf$ on $\R$.  Equivalently, the associated finite
probability measures converge weakly to the probability measure with density
$\upf$.
\end{theorem}

\begin{proof}
Normalize $p_n$ by $2^{-\binom{n+1}{2}}$.  Its coefficients sum to $1$, so
replacing $X^m$ by a point mass at
\[
 \frac{2m-g_n}{2^{n+1}}
\]
produces a probability measure $\mu_n$.  Formula \eqref{eq:p-products}
identifies this measure with the law of a finite sum of independent symmetric
lattice variables.  Its characteristic function is the corresponding finite
partial product in \eqref{eq:char-product}; hence $\mu_n$ converges weakly to
the law with density $\upf$.

The step function \eqref{eq:step-approx} is obtained by spreading every atom
uniformly over a cell of width $2^{-n}$.  Integrals of $\upf_n$ against
indicators of intervals with dyadic endpoints therefore agree with the
corresponding masses up to a boundary term that tends to zero.  Weak
convergence gives convergence of these integrals to those of $\upf$.
Symmetry and the log-concave/unimodal coefficient recurrence inherited from
\eqref{eq:p-recurrence} show that $\upf_n$ is nondecreasing on $(-1,0)$ and
nonincreasing on $(0,1)$.  Since the limit is continuous, convergence of the
interval integrals plus monotonicity gives pointwise convergence.  The
normalization at zero follows directly from the central coefficient
recurrence.
\end{proof}

\begin{remark}
The step approximation is not merely a picture-making device.  It gives a
finite combinatorial model for $\upf$, makes positivity visible before any
Fourier analysis, and is particularly convenient for certified numerical
bounds because all breakpoints and heights are dyadic rationals times
integers.
\end{remark}

\subsection{Partition of unity and Poisson summation}

The rapid decay of $\widehat{\upf}$ permits Poisson summation without any
regularization.

\begin{theorem}[Poisson formula]\label{thm:poisson-up}
For $u>0$ and $t\in\R$,
\begin{equation}\label{eq:poisson-up}
 \sum_{k\in\Z}\upf(t+uk)
 =\frac1u\sum_{m\in\Z}
   \widehat{\upf}(m/u)\e^{2\pi\ii mt/u}.
\end{equation}
Both series converge absolutely after the usual interpretation of the
locally finite left side.
\end{theorem}

\begin{proof}
Apply the classical Poisson formula to the Schwartz function $\upf$.  Although
$\upf$ is compactly supported rather than rapidly decaying in the usual
nontrivial way, it is still a Schwartz function because every derivative is
compactly supported.  Our Fourier convention gives exactly the phase in
\eqref{eq:poisson-up}.
\end{proof}

\begin{proposition}[Dyadic and integral partitions of unity]
\label{prop:partition-unity}
For every positive integer $n$,
\begin{equation}\label{eq:fine-partition}
 \sum_{k\in\Z}\upf\left(t+\frac{k}{n}\right)=n.
\end{equation}
In particular,
\begin{equation}\label{eq:unit-partition}
 \sum_{k\in\Z}\upf(t+k)=1,
 \qquad
 \upf(t)+\upf(t-1)=1\quad(0\le t\le1).
\end{equation}
\end{proposition}

\begin{proof}
Take $u=1/n$ in \eqref{eq:poisson-up}.  For every nonzero integer $m$,
$\widehat{\upf}(nm)=0$, because the first sinc factor in
\eqref{eq:char-product} vanishes.  Only the zero Fourier mode remains and has
value $n\widehat{\upf}(0)=n$.  The support restriction reduces the case $n=1$
to the two-term identity on $[0,1]$.
\end{proof}

The last identity is another direct proof that
$F(x)=\upf(x-1)$ satisfies $F(1-x)=1-F(x)$.

\subsection{Cosine and Weierstrass products}

The elementary identity
\[
 \frac{\sin z}{z}=\prod_{j=1}^{\infty}\cos(z/2^j)
\]
and a rearrangement of absolutely convergent products give
\begin{equation}\label{eq:cos-product}
 \widehat{\upf}(z)
 =\prod_{m=1}^{\infty}\cos(\pi z/2^m)^m.
\end{equation}
Combining this with Euler's product for the sine gives a second factorization:
\begin{equation}\label{eq:weierstrass-product}
 \widehat{\upf}(z)
 =\prod_{r=1}^{\infty}
   \left(1-\frac{z^2}{r^2}\right)^{1+\nu_2(r)},
\end{equation}
where $\nu_2(r)$ is the exponent of $2$ in $r$.  Indeed, a fixed integer
$r$ occurs once for every representation $r=2^h m$, namely
$0\le h\le\nu_2(r)$.

Taking $u=2$ in Poisson summation, using support and pairing conjugate modes,
we obtain the uniformly convergent cosine expansion on $[-1,1]$:
\begin{equation}\label{eq:cosine-series-up}
 \upf(t)=\frac12+\sum_{m=0}^{\infty}
 \widehat{\upf}\left(\frac{2m+1}{2}\right)
 \cos\bigl((2m+1)\pi t\bigr).
\end{equation}
All nonzero even Fourier samples disappear because they are integer zeros of
$\widehat{\upf}$.  Formula \eqref{eq:cosine-series-up} is useful for high
precision evaluation away from the endpoints; exact dyadic evaluation is
better handled arithmetically in \Cref{sec:dyadic-closed-form}.

\section{The global extension and the derivative tower}
\label{sec:global-extension}

We return to the locally finite series
\begin{equation}\label{eq:theta-def}
 \Th(x)=\sum_{m\ge0}\epsilon_m\upf(x-2m-1),
 \qquad \epsilon_m=(-1)^{\wt(m)}.
\end{equation}
The binary identities
\begin{equation}\label{eq:tm-binary}
 \epsilon_{2m}=\epsilon_m,
 \qquad \epsilon_{2m+1}=-\epsilon_m
\end{equation}
are exactly what is needed to turn the two translates in the Rvachev equation
into one scaled copy of $\Th$.

\begin{theorem}[Global scaling equation]\label{thm:theta-scaling}
The function $\Th$ is $C^\infty$ and
\begin{equation}\label{eq:theta-prime}
 \Th'(x)=2\Th(2x),\qquad x\in\R.
\end{equation}
For every $k\ge0$,
\begin{equation}\label{eq:theta-derivatives}
 \Th^{(k)}(x)=2^{\binom{k+1}{2}}\Th(2^kx).
\end{equation}
Moreover, for $-1\le t\le1$,
\begin{equation}\label{eq:up-derivatives-theta}
 \upf^{(k)}(t)=2^{\binom{k+1}{2}}
 \Th\bigl(2^kt+2^k\bigr).
\end{equation}
\end{theorem}

\begin{proof}
Local finiteness permits termwise differentiation.  Using
\eqref{eq:up-diffeq},
\begin{align*}
 \Th'(x)
 &=2\sum_{m\ge0}\epsilon_m
 \bigl[\upf(2x-4m-1)-\upf(2x-4m-3)\bigr]\\
 &=2\sum_{m\ge0}\epsilon_{2m}\upf(2x-2(2m)-1)
   +2\sum_{m\ge0}\epsilon_{2m+1}\upf(2x-2(2m+1)-1)\\
 &=2\Th(2x).
\end{align*}
Induction gives \eqref{eq:theta-derivatives}, since differentiating the
$k$th identity contributes an additional factor $2^{k+1}$.  Finally,
$\upf(t)=\Th(t+1)$ on $[-1,1]$; differentiating and applying
\eqref{eq:theta-derivatives} at $t+1$ proves
\eqref{eq:up-derivatives-theta}.
\end{proof}

At integers, the support of $\upf$ makes the global extension especially
simple:
\begin{equation}\label{eq:theta-integers}
 \Th(2m)=0,
 \qquad
 \Th(2m+1)=\epsilon_m
 \quad(m\ge0).
\end{equation}
These values explain both the terminating Taylor expansions at dyadic points
and the exact-arithmetic formulas later in the article.

\begin{corollary}[Terminating Taylor data at dyadic points]
\label{cor:dyadic-taylor}
Let $t=q/2^n\in(-1,1)$ be in lowest dyadic terms, so $q$ is odd.  Then
\begin{equation}\label{eq:dyadic-derivative-end}
 \upf^{(k)}(t)=0\quad(k>n),
 \qquad
 \upf^{(n)}(t)\ne0.
\end{equation}
Consequently the formal Taylor series of $\upf$ at $t$ is the polynomial
\begin{equation}\label{eq:dyadic-taylor-poly}
 T_t(h)=\sum_{k=0}^{n}\frac{\upf^{(k)}(t)}{k!}h^k.
\end{equation}
All Taylor coefficients after degree $n$ vanish.
\end{corollary}

\begin{proof}
For $k\ge n$, the argument of $\Th$ in
\eqref{eq:up-derivatives-theta} is
\[
 2^kt+2^k=2^{k-n}q+2^k.
\]
At $k=n$ this is the odd integer $q+2^n$, and
\eqref{eq:theta-integers} gives a nonzero value.  For $k>n$ it is even, so
the value is zero.
\end{proof}

\begin{theorem}[Nowhere analyticity]\label{thm:nowhere-analytic}
The function $\upf$ is not real-analytic at any point of $[-1,1]$.  Likewise,
$F$ is not real-analytic at any point of $[0,1]$.
\end{theorem}

\begin{proof}
Suppose $\upf$ were analytic on a neighborhood of some point.  That
neighborhood contains a reduced dyadic point $q/2^n$.  By
\Cref{cor:dyadic-taylor}, analyticity there would force $\upf$ to agree on a
smaller interval with a polynomial of degree at most $n$.  The same interval
contains reduced dyadic points of arbitrarily large denominator exponent
$N$.  At such a point the $N$th derivative is nonzero by
\eqref{eq:dyadic-derivative-end}, whereas every derivative of order greater
than $n$ of the polynomial vanishes.  This contradiction proves the claim.
The assertion for $F$ follows by translation on each half of the unit
interval.
\end{proof}

\begin{remark}[Flat but not analytic]
At the boundary, $F^{(k)}(0)=0$ for every $k\ge0$, yet $F(x)>0$ for every
$x>0$.  Thus the Taylor series at zero is identically zero.  The endpoint
asymptotic in \Cref{sec:asymptotics} quantifies this flatness much more
precisely: it is essentially the exponential of a negative quadratic in
$\log(1/x)$, modified by logarithmic and periodic terms.
\end{remark}

\section{Moments, generating functions, and inverse powers of two}
\label{sec:moments}

\subsection{Even moments}

Define
\begin{equation}\label{eq:c-def}
 c_n=\int_{-1}^{1}t^{2n}\upf(t)\dd t,
 \qquad n\ge0.
\end{equation}
Then $c_0=1$, and the Taylor series of the Fourier transform is
\begin{equation}\label{eq:fourier-moment-series}
 \widehat{\upf}(z)=
 \sum_{n=0}^{\infty}(-1)^n c_n\frac{(2\pi z)^{2n}}{(2n)!}.
\end{equation}

\begin{proposition}[Moment recurrence]\label{prop:c-recurrence}
For $n\ge1$,
\begin{equation}\label{eq:c-recurrence}
 (2n+1)(2^{2n}-1)c_n
 =\sum_{k=0}^{n-1}\binom{2n+1}{2k}c_k.
\end{equation}
The first values are
\begin{equation}\label{eq:c-first}
 1,\quad \frac19,\quad \frac{19}{675},\quad
 \frac{583}{59535},\quad
 \frac{132809}{32531625},\quad
 \frac{46840699}{24405225075},\ldots
\end{equation}
\end{proposition}

\begin{proof}
Insert \eqref{eq:fourier-moment-series} into the refinement equation
\eqref{eq:fourier-refinement}, expand
$\sin(\pi z)/(\pi z)$, and compare the coefficient of $z^{2n}$.  Equivalently,
coefficient comparison first gives
\[
 (2n+1)2^{2n}c_n
 =\sum_{k=0}^{n}\binom{2n+1}{2k}c_k.
\]
The final summand is $(2n+1)c_n$; moving it to the left yields
\eqref{eq:c-recurrence}.
\end{proof}

Since $\upf$ is even,
\begin{equation}\label{eq:half-even-moment}
 \int_0^1t^{2n}\upf(t)\dd t=\frac{c_n}{2}.
\end{equation}

\subsection{Half-moments and a generating function}

Define $d_0=1$ and, for $n\ge1$,
\begin{equation}\label{eq:d-integral}
 d_n=n\int_0^1t^{n-1}\upf(t)\dd t.
\end{equation}
Introduce
\begin{equation}\label{eq:f-generating}
 f(z)=1+z\int_0^1\e^{zt}\upf(t)\dd t
     =\sum_{n=0}^{\infty}d_n\frac{z^n}{n!}.
\end{equation}
The probability interpretation gives another form:
\begin{equation}\label{eq:f-fourier-form}
 f(z)=\e^{z/2}\widehat{\upf}\left(\frac{\ii z}{4\pi}\right).
\end{equation}
Indeed, the right side is the moment-generating function of the random
variable $X$ in \eqref{eq:X-random-series}.  Using
$\upf(t)=1-F(t)$ on $[0,1]$ and integration by parts gives the first
expression in \eqref{eq:f-generating}.

\begin{proposition}[Refinement of the generating function]
\label{prop:f-refinement}
The entire function $f$ satisfies
\begin{equation}\label{eq:f-refinement}
 f(2z)=\frac{\e^z-1}{z}f(z),
\end{equation}
with the quotient understood by continuity at $z=0$.  Therefore
\begin{equation}\label{eq:d-recurrence}
 d_0=1,
 \qquad
 (n+1)(2^n-1)d_n
 =\sum_{k=0}^{n-1}\binom{n+1}{k}d_k
 \quad(n\ge1).
\end{equation}
\end{proposition}

\begin{proof}
Substitute \eqref{eq:f-fourier-form} into
\eqref{eq:fourier-refinement} at the imaginary argument $\ii z/(2\pi)$.
The hyperbolic-sine simplification gives \eqref{eq:f-refinement}.  Expanding
both sides in exponential generating series and comparing coefficients gives
\eqref{eq:d-recurrence}.
\end{proof}

The first values are
\begin{equation}\label{eq:d-first}
 1,\quad\frac12,\quad\frac5{18},\quad\frac16,\quad
 \frac{143}{1350},\quad\frac{19}{270},\quad
 \frac{1153}{23814},\quad\frac{583}{17010},\ldots
\end{equation}

\begin{theorem}[Moment--endpoint identity]\label{thm:moment-endpoint}
For every $n\ge1$,
\begin{equation}\label{eq:moment-endpoint}
 \int_0^1t^{n-1}\upf(t)\dd t
 =(n-1)!\,2^{\binom n2}F(2^{-n}).
\end{equation}
Equivalently,
\begin{equation}\label{eq:d-F-inverse}
 d_n=n!\,2^{\binom n2}F(2^{-n})
    =n!\,2^{\binom n2}\upf(1-2^{-n}).
\end{equation}
\end{theorem}

\begin{proof}
All derivatives of $F$ vanish at $0$.  Taylor's theorem with integral
remainder, applied at order $n-1$ and at $a=2^{-n}$, gives
\[
 F(2^{-n})=\frac1{(n-1)!}
 \int_0^{2^{-n}}(2^{-n}-t)^{n-1}F^{(n)}(t)\dd t.
\]
Repeated differentiation of $F'(t)=2F(2t)$ is valid throughout
$0\le t\le2^{-n}$ and gives
\[
 F^{(n)}(t)=2^{\binom{n+1}{2}}F(2^nt).
\]
Set $u=2^nt$.  The powers of $2$ simplify to
$2^{-\binom n2}$, and hence
\[
 F(2^{-n})=\frac{2^{-\binom n2}}{(n-1)!}
 \int_0^1(1-u)^{n-1}F(u)\dd u.
\]
Finally $F(u)=\upf(u-1)=\upf(1-u)$, and the substitution $t=1-u$ proves
\eqref{eq:moment-endpoint}.  Multiplication by $n$ gives
\eqref{eq:d-F-inverse}.
\end{proof}

\begin{proposition}[Half-moments from even moments]\label{prop:d-from-c}
For every $n\ge0$,
\begin{equation}\label{eq:d-from-c}
 d_n=2^{-n}\sum_{k=0}^{\lfloor n/2\rfloor}
 \binom{n}{2k}c_k.
\end{equation}
\end{proposition}

\begin{proof}
For $n\ge1$, integrate \eqref{eq:d-integral} by parts.  The boundary term
vanishes because $\upf(1)=0$, and the Rvachev equation on $(0,1)$ gives
\begin{align*}
 d_n
 &=-\int_0^1t^n\upf'(t)\dd t
   =2\int_0^1t^n\upf(2t-1)\dd t\\
 &=\int_{-1}^{1}\left(\frac{x+1}{2}\right)^n\upf(x)\dd x.
\end{align*}
Expand $(x+1)^n$.  Odd moments vanish by evenness, and the even moments are
$c_k$.  The case $n=0$ is immediate.
\end{proof}

\section{A closed formula for every dyadic value}
\label{sec:dyadic-closed-form}

The rationality of dyadic values is often proved recursively.  A more
illuminating formula evaluates them directly from the even moments $c_k$ and
a finite Thue--Morse sum.  It is valid not only on the unit interval but for
the signed global extension.

\begin{definition}\label{def:Qna}
For $n,a\in\N$, define
\begin{equation}\label{eq:Qna}
 Q_n(a)=\frac{2^{-\binom{n+1}{2}}}{n!}
 \sum_{h=0}^{a-1}\epsilon_h
 \sum_{k=0}^{\lfloor n/2\rfloor}
 \binom{n}{2k}(2a-2h-1)^{n-2k}c_k,
\end{equation}
where the outer sum is empty when $a=0$.
\end{definition}

\begin{theorem}[Exact dyadic formula]\label{thm:closed-dyadic}
For every $n,a\in\N$,
\begin{equation}\label{eq:theta-Q}
 \Th\left(\frac{a}{2^n}\right)=Q_n(a).
\end{equation}
Consequently, if $0\le a\le2^n$, then
\begin{equation}\label{eq:F-Q}
 F\left(\frac{a}{2^n}\right)=Q_n(a)\in\Q.
\end{equation}
For an integer $q$,
\begin{equation}\label{eq:up-Q}
 \upf\left(\frac{q}{2^n}\right)=
 \begin{cases}
 Q_n(2^n-|q|),&|q|\le2^n,\\
 0,&|q|>2^n.
 \end{cases}
\end{equation}
\end{theorem}

\begin{proof}
Because $\Th$ and all its derivatives vanish on the negative half-line,
repeated integration of its $(n+1)$st derivative gives, for every real $x$,
\begin{equation}\label{eq:theta-repeated-integral}
 \Th(x)=\frac1{n!}\int_{-\infty}^{x}(x-t)^n\Th^{(n+1)}(t)\dd t.
\end{equation}
By \eqref{eq:theta-derivatives},
\[
 \Th^{(n+1)}(t)=2^{\binom{n+2}{2}}\Th(2^{n+1}t).
\]
Set $x=a/2^n$ and then $u=2^{n+1}t$.  The powers of $2$ cancel to
$2^{-\binom{n+1}{2}}$, yielding
\begin{equation}\label{eq:theta-integral-a}
 \Th\left(\frac{a}{2^n}\right)
 =\frac{2^{-\binom{n+1}{2}}}{n!}
  \int_{-\infty}^{2a}(2a-u)^n\Th(u)\dd u.
\end{equation}
Insert the locally finite bump decomposition \eqref{eq:theta-def}.  Up to the
upper limit $2a$, exactly the complete bumps indexed by $0\le h<a$ occur:
\begin{align*}
 \int_{-\infty}^{2a}(2a-u)^n\Th(u)\dd u
 &=\sum_{h=0}^{a-1}\epsilon_h
   \int_{2h}^{2h+2}(2a-u)^n\upf(u-2h-1)\dd u.
\end{align*}
With $v=u-2h-1$, each integral becomes
\[
 \int_{-1}^{1}(2a-2h-1-v)^n\upf(v)\dd v.
\]
Expand by the binomial theorem.  The odd moments vanish by evenness, while
\[
 \int_{-1}^{1}v^{2k}\upf(v)\dd v=c_k.
\]
The result is precisely the double sum in \eqref{eq:Qna}.

On $[0,1]$, $\Th=F$, proving \eqref{eq:F-Q}.  Finally, evenness and the
relation $\upf(t)=F(1-|t|)$ for $|t|\le1$ give
\eqref{eq:up-Q}.
\end{proof}

\begin{remark}[Why the formula is finite]
The finiteness has two independent sources.  Only $a$ translated bumps lie
to the left of $2a$, and only even moments up to order $n$ appear in the
binomial expansion.  No limiting operation and no numerical approximation
remains.  Since the recurrence \eqref{eq:c-recurrence} computes every $c_k$
in $\Q$, the formula is an exact rational algorithm.
\end{remark}

\subsection{Examples}

The first two nontrivial levels are
\begin{align*}
 F(1/2)&=\frac12,\\
 F(1/4)&=\frac5{72},&F(3/4)&=\frac{67}{72},\\
 F(1/8)&=\frac1{288},&F(3/8)&=\frac{73}{288},
 &F(5/8)&=\frac{215}{288},&F(7/8)&=\frac{287}{288}.
\end{align*}
At denominator $16$, selected values are
\begin{center}
\begin{tabular}{@{}c@{\qquad}c@{\qquad}c@{\qquad}c@{}}
\toprule
$x$ & $F(x)$ & $x$ & $F(x)$\\
\midrule
$1/16$ & $143/2073600$ & $9/16$ & $1295857/2073600$\\
$2/16$ & $1/288$       & $10/16$& $215/288$\\
$3/16$ & $46657/2073600$ & $11/16$& $1767743/2073600$\\
$4/16$ & $5/72$        & $12/16$& $67/72$\\
$5/16$ & $305857/2073600$ & $13/16$& $2026943/2073600$\\
$6/16$ & $73/288$      & $14/16$& $287/288$\\
$7/16$ & $777743/2073600$ & $15/16$& $2073457/2073600$\\
$8/16$ & $1/2$         & $16/16$& $1$\\
\bottomrule
\end{tabular}
\end{center}
The table exhibits both refinement invariance, for example
$F(2/16)=F(1/8)$, and reflection $F(1-x)=1-F(x)$.

\subsection{Refinement invariance}

A dyadic rational has many presentations.  Correctness requires
\begin{equation}\label{eq:Q-refinement}
 Q_{n+1}(2a)=Q_n(a).
\end{equation}
This identity follows immediately from \Cref{thm:closed-dyadic}, but it can
also be proved purely algebraically from the Thue--Morse relations
\eqref{eq:tm-binary} and the moment recurrence.  In the formal development,
both the closed expression and the executable evaluator are proved invariant
under arbitrary repeated refinements $(a,n)\mapsto(2a,n+1)$; thus exact
computation does not depend on first reducing the dyadic fraction.

\section{A fast bit-recursive exact evaluator}\label{sec:bit-recursion}

The closed formula sums over every $h<a$.  It is excellent for theory but can
be wasteful when $a$ has only a few nonzero binary digits.  A terminating
Taylor identity removes one highest set bit at each step.

\subsection{Taylor data at inverse powers of two}

For $0\le k\le r$, repeated differentiation on the left half and
\eqref{eq:d-F-inverse} give
\begin{equation}\label{eq:derivative-inverse-power}
 F^{(k)}(2^{-r})
 =2^{\binom{k+1}{2}}F(2^{-(r-k)})
 =2^{\binom{k+1}{2}-\binom{r-k}{2}}
   \frac{d_{r-k}}{(r-k)!}.
\end{equation}
All derivatives of order greater than $r$ vanish at $2^{-r}$.

\begin{theorem}[Terminating Taylor recursion]\label{thm:taylor-recursion}
Let $r\ge1$ and $0\le y<2^{-r}$.  Then
\begin{equation}\label{eq:taylor-recursion}
 F(2^{-r}+y)
 =-F(y)+\sum_{k=0}^{r}
 2^{\binom{k+1}{2}-\binom{r-k}{2}}
 \frac{d_{r-k}}{(r-k)!}\frac{y^k}{k!}.
\end{equation}
More generally, the same formula holds for the signed global extension on the
corresponding dyadic block, with the empty sum convention in the ranges where
the Taylor center is an even integer.
\end{theorem}

\begin{proof}
Apply Taylor's theorem with integral remainder of order $r$ at $2^{-r}$.
The polynomial part is \eqref{eq:derivative-inverse-power}.  It remains to
identify the remainder.  Using $F=\Th$ on the unit interval and
\eqref{eq:theta-derivatives}, the $(r+1)$st derivative on the short interval
to the right of $2^{-r}$ is a rescaled copy of the next signed bump.  More
explicitly, after the change of variables dictated by
$\Th^{(r+1)}(t)=2^{\binom{r+2}{2}}\Th(2^{r+1}t)$, the integral remainder is
\[
 -\frac1{r!}\int_0^y(y-v)^r F^{(r+1)}(v)\dd v.
\]
Taylor's integral formula at the flat point $0$ identifies the last integral
with $-F(y)$.  This is the analytic origin of the minus sign in
\eqref{eq:taylor-recursion}.
\end{proof}

\subsection{Precomputation at inverse powers}

Let
\[
 V_j=F(2^{-j}),\qquad V_0=1.
\]
The $d$-recurrence or \eqref{eq:d-F-inverse} computes these values.  A direct
triangular recurrence used by the exact evaluator is
\begin{equation}\label{eq:V-recurrence}
 V_{n+1}=\frac{1}{2^{n+1}-1}
 \sum_{k=0}^{n}
 \frac{V_k}{
 2^{\binom{n+1}{2}-\binom{k}{2}}(n+2-k)!}.
\end{equation}
For example, $V_1=1/2$, $V_2=5/72$, $V_3=1/288$, and
$V_4=143/2073600$.

\subsection{The binary algorithm}

Suppose $0<a<2^n$ and let $2^\ell$ be the highest power of two not exceeding
$a$.  Put
\[
 r=n-\ell,
 \qquad b=a-2^\ell,
 \qquad y=\frac{b}{2^n}.
\]
Then $a/2^n=2^{-r}+y$ and $0\le y<2^{-r}$, so
\eqref{eq:taylor-recursion} reduces $F(a/2^n)$ to a known Taylor polynomial
minus $F(b/2^n)$.  The numerator $b$ is obtained by deleting the highest set
bit of $a$.

\begin{algorithm}[Exact bounded dyadic evaluation]\label{alg:dyadic-bit}
To compute $F(a/2^n)$ exactly:
\begin{enumerate}
\item Return $0$ if $a\le0$ and $1$ if $a\ge2^n$.
\item Precompute $V_0,\ldots,V_n$ by \eqref{eq:V-recurrence}, or equivalently
      precompute $d_0,\ldots,d_n$ by \eqref{eq:d-recurrence}.
\item If $0<a<2^n$, remove the highest set bit as above, evaluate the finite
      Taylor polynomial in \eqref{eq:taylor-recursion} by Horner's rule, and
      subtract the recursive value at $b/2^n$.
\end{enumerate}
The recursion depth is exactly the number of set bits of $a$.
\end{algorithm}

Thus a numerator such as $a=2^{100}+2^7+1$ requires only three recursive
Taylor steps, whereas the closed sum \eqref{eq:Qna} contains roughly
$2^{100}$ outer terms if used naively.

\subsection{Bounded, global, and rational interfaces}

For the bounded function, signed dyadic arguments are clamped according to
Definition \ref{def:bounded-fabius}.  The signed global extension has a
different right-hand behavior.  Let $s=2^n$, let $a>0$, and write
\[
 a=2sB+R,
 \qquad 0\le R<2s.
\]
Set $C=R$ if $R\le s$ and $C=2s-R$ otherwise.  Then
\begin{equation}\label{eq:global-triangular-reduction}
 \Th(a/2^n)=\epsilon_B F(C/2^n).
\end{equation}
This is the triangular reduction modulo $2$, followed by the Thue--Morse sign
of the block.  It gives a total exact evaluator for $\Th$ on all dyadic
arguments.

A reduced rational $p/q$ is dyadic exactly when $q$ is a power of two.  Hence
an exact rational interface first checks this denominator condition, returns
``not dyadic'' otherwise, and invokes the signed-integer dyadic evaluator when
the condition holds.  No floating-point recognition of powers of two is
needed.

\section{The finite \texorpdfstring{$q$}{q}-binomial--Thue--Morse formula}
\label{sec:qbinomial}

The formula of \Cref{thm:closed-dyadic} uses ordinary binomial coefficients
and moments.  A different expression, motivated by a conjectured symbolic
formula, uses Gaussian binomial coefficients at $q=1/2$ and power sums with
Thue--Morse signs.  Its proof in the formal development is stronger than the
original statement: an arbitrary common translation may be inserted in all
inner powers without changing the answer.

For $q\in\Q$, define
\begin{equation}\label{eq:qpochhammer}
 (q;q)_n=\prod_{j=1}^{n}(1-q^j),
 \qquad
 \genfrac{[}{]}{0pt}{}{n}{k}_q=\frac{(q;q)_n}{(q;q)_k(q;q)_{n-k}}
 \quad(0\le k\le n).
\end{equation}
For $c\in\Q$ and natural $m,k,p$, put
\begin{equation}\label{eq:translated-power-sum}
 S_{m,k,p}(c)=
 \sum_{j=0}^{m2^k-1}\epsilon_j
 \bigl(j-m2^k+c\bigr)^p.
\end{equation}

\begin{theorem}[$q$-binomial formula]\label{thm:qbinomial-formula}
For all $m,n\in\N$ and every common shift $c\in\Q$,
\begin{equation}\label{eq:qbinomial-formula}
 \Th\left(\frac{m}{2^n}\right)
 =\frac{1}{2^{n^2}(1/2;1/2)_n}
 \sum_{k=0}^{n}
 \frac{\qbinom{n}{k}}
      {2^{k(k-1)}(n+k)!}
 S_{m,k,n+k}(c).
\end{equation}
The right side is independent of $c$.  If $m\le2^n$, it is the bounded value
$F(m/2^n)$.  In particular, $m=1$ gives a finite formula for
$F(2^{-n})$.
\end{theorem}

\begin{proof}
We give the algebraic mechanism in several steps.

\emph{Step 1: exponential generating functions for the inner sums.}
For fixed $m$ and $k$,
\begin{equation}\label{eq:S-egf}
 \sum_{p\ge0}S_{m,k,p}(c)\frac{z^p}{p!}
 =\e^{(c-m2^k)z}
  \sum_{j=0}^{m2^k-1}\epsilon_j\e^{jz}.
\end{equation}
Splitting $j=h2^k+r$ and using
$\epsilon_{h2^k+r}=\epsilon_h\epsilon_r$ gives the block factorization
\begin{equation}\label{eq:tm-block-factor}
 \sum_{j=0}^{m2^k-1}\epsilon_j\e^{jz}
 =\left(\sum_{h=0}^{m-1}\epsilon_h\e^{h2^kz}\right)
  \prod_{r=0}^{k-1}(1-\e^{2^rz}).
\end{equation}
The product on the right has a zero of exact order $k$ at $z=0$; this is the
classical vanishing of Thue--Morse power sums below degree $k$.

\emph{Step 2: remove the forced zero.}
Divide the product in \eqref{eq:tm-block-factor} by $z^k$.  Its constant term
is $(-1)^k2^{\binom k2}$.  After this normalization, the coefficient of
$z^n$ is exactly
$S_{m,k,n+k}(c)/(n+k)!$ up to the displayed power of $2$ and sign.  Thus the
outer sum in \eqref{eq:qbinomial-formula} is a weighted coefficient extraction
from a family of refinement-factor series indexed by $k$.

\emph{Step 3: the Gaussian-binomial annihilator.}
The finite $q$-binomial theorem
\begin{equation}\label{eq:finite-q-binomial}
 \sum_{k=0}^{n}(-1)^kq^{\binom k2}\genfrac{[}{]}{0pt}{}{n}{k}_qX^k
 =\prod_{j=0}^{n-1}(1-Xq^j)
\end{equation}
shows that the weights at the nodes $2^{-k}$ annihilate every polynomial in
that node of degree less than $n$ and evaluate the degree-$n$ term by
\begin{equation}\label{eq:qweight-self}
 2^{\binom{n+1}{2}}(1/2;1/2)_n.
\end{equation}
Applied coefficientwise to the normalized refinement factors, all lower
coefficients disappear.  The surviving coefficient is the $n$th coefficient
of the product of two series: the finite Thue--Morse prefix series associated
with $m$, and the recurrence series
\[
 \sum_{r\ge0}d_r\frac{z^r}{r!},
\]
which satisfies the refinement equation \eqref{eq:f-refinement}.

\emph{Step 4: identify the surviving coefficient.}
The convolution coefficient just described is the exact global dyadic value
$\Th(m/2^n)$ in the recurrence normalization.  The scalar in
\eqref{eq:qweight-self}, together with the powers extracted in Step 2,
simplifies to the denominator
$2^{n^2}(1/2;1/2)_n$ in
\eqref{eq:qbinomial-formula}.

\emph{Step 5: translation invariance.}
Replacing $c=0$ by a common $c$ multiplies every inner generating function by
$\e^{cz}$.  In the coefficient convolution, the terms involving a positive
power of $z$ from $\e^{cz}$ are paired with weighted coefficients of degree
strictly below $n$, and these are annihilated by Step 3.  Only the constant
term $1$ survives.  Hence the complete expression is independent of $c$.
This proves the theorem.
\end{proof}

\begin{remark}
The source formula uses $c=1/2$, giving powers
$(j-m2^k+1/2)^{n+k}$.  The centered choice $c=0$ is often simpler for parity
arguments.  The theorem says that this is not a heuristic simplification but
an exact invariance of the full outer sum.  The formal proof also transports
the identity from $\Q$ to $\R$ and $\C$ by coercion, after the rational
identity has been established.
\end{remark}

\section{Arithmetic sequences, denominators, and parity}
\label{sec:arithmetic}

\subsection{Integer normalizations of the moments}

Let $(2n+1)!!=1\cdot3\cdots(2n+1)$.  Define
\begin{equation}\label{eq:Fn-def}
 \mathcal F_n
 =c_n(2n+1)!!\prod_{\nu=1}^{n}(2^{2\nu}-1).
\end{equation}
We use the calligraphic symbol $\mathcal F_n$ to avoid confusion with the
function $F$.

\begin{proposition}\label{prop:Fn-integer}
Every $\mathcal F_n$ is a positive integer.  More precisely,
$\mathcal F_0=1$ and, for $n\ge1$,
\begin{equation}\label{eq:Fn-recurrence}
 \mathcal F_n=
 \sum_{k=0}^{n-1}\mathcal F_k\binom{2n+1}{2k}
 \frac{(2n-1)!!}{(2k+1)!!}
 \prod_{\nu=k+1}^{n-1}(2^{2\nu}-1).
\end{equation}
\end{proposition}

\begin{proof}
Substitute \eqref{eq:Fn-def} into the moment recurrence
\eqref{eq:c-recurrence} and clear the common factors.  Every factor on the
right of \eqref{eq:Fn-recurrence} is an integer; induction proves the claim.
\end{proof}

Similarly define $G_n$ by
\begin{equation}\label{eq:Gn-def}
 d_n=\frac{G_n}{(n+1)!\prod_{j=1}^{n}(2^j-1)}.
\end{equation}
Then $G_0=1$ and
\begin{equation}\label{eq:Gn-recurrence}
 G_n=\sum_{k=0}^{n-1}G_k\binom{n+1}{k}
 \frac{n!}{(k+1)!}
 \prod_{j=k+1}^{n-1}(2^j-1),
\end{equation}
so all $G_n$ are positive integers.  The first terms are
\[
 1,\ 1,\ 5,\ 84,\ 4004,\ 494760,\ 150120600,
 \ 107969547840,\ldots
\]
Combining \eqref{eq:Gn-def} with \eqref{eq:d-F-inverse} gives an exact
integer-normalized formula for the inverse-power values:
\begin{equation}\label{eq:F-inverse-G}
 F(2^{-n})=
 \frac{2^{-\binom n2}G_n}
 {n!(n+1)!\prod_{j=1}^{n}(2^j-1)}.
\end{equation}

\subsection{Reshetnikov's integers}

For $n\ge1$, define
\begin{equation}\label{eq:Rn-def}
 R_n=2^{\binom{n-1}{2}}(2n)!F(2^{-n})
 \prod_{j=1}^{\lfloor n/2\rfloor}(2^{2j}-1).
\end{equation}
The positive exponent in \eqref{eq:Rn-def} is essential; a sign typo in one
summary of the formula reverses it.

Using \eqref{eq:d-F-inverse} and
$(2n)!=2^n n!(2n-1)!!$ gives
\begin{equation}\label{eq:Rn-dn}
 R_n=2d_n(2n-1)!!
 \prod_{j=1}^{\lfloor n/2\rfloor}(2^{2j}-1).
\end{equation}
For odd indices,
\begin{equation}\label{eq:R-odd}
 R_{2n+1}=\mathcal F_n\frac{(4n+1)!!}{(2n-1)!!}.
\end{equation}
For even indices,
\begin{equation}\label{eq:R-even}
 R_{2n}=\sum_{k=0}^{n}
 \frac{2\mathcal F_k}{2^{2n}}\binom{2n}{2k}
 \frac{(4n-1)!!}{(2k+1)!!}
 \prod_{\ell=k+1}^{n}(2^{2\ell}-1).
\end{equation}
The latter formula shows first that the only possible denominator is a power
of two.

\begin{theorem}[Integrality of $R_n$]\label{thm:R-integer}
Every $R_n$ is a positive integer.
\end{theorem}

\begin{proof}
The odd case follows at once from \eqref{eq:R-odd}.  By
\eqref{eq:Rn-dn}, $\nu_2(R_n)=\nu_2(2d_n)$ because the remaining factors are
odd.  We prove $\nu_2(2d_n)\ge0$ by induction.  Odd indices are already
covered by \eqref{eq:R-odd}, and $2d_0=2$.  For an even index $2n$, multiply
\eqref{eq:d-recurrence} by $2$ to obtain
\begin{equation}\label{eq:2d-even-induction}
 (2n+1)(2^{2n}-1)2d_{2n}
 =\sum_{k=0}^{2n-1}\binom{2n+1}{k}2d_k.
\end{equation}
The left prefactor is odd.  By induction every summand on the right is
$2$-integral, hence so is $2d_{2n}$.  Formula \eqref{eq:R-even} says there are
no odd-prime denominators, so $R_{2n}$ is an integer.  Positivity follows from
$F(2^{-n})>0$.
\end{proof}

\subsection{Denominator bounds for all dyadic values}

Let $D_n$ be the least common multiple of the reduced denominators of
$\upf(a/2^n)$ for odd integers $a$ with $|a|<2^n$.  Refinement invariance
shows that this captures the genuinely new denominator at level $n$.

\begin{theorem}[Uniform denominator bound]\label{thm:denominator-bound}
For $n\ge1$,
\begin{equation}\label{eq:D-bound}
 D_n\mid
 n!\,2^{\binom{n+1}{2}}
 (2\lfloor n/2\rfloor+1)!!
 \prod_{k=1}^{\lfloor n/2\rfloor}(2^{2k}-1).
\end{equation}
Moreover, $D_n$ is a multiple of the denominator of
\begin{equation}\label{eq:D-lower}
 \frac{d_n}{2^{\binom n2}n!}=F(2^{-n}).
\end{equation}
\end{theorem}

\begin{proof}
After multiplication by $n!2^{\binom{n+1}{2}}$, the closed formula
\eqref{eq:Qna} is an integral linear combination of
$c_0,\ldots,c_{\lfloor n/2\rfloor}$.  Formula \eqref{eq:Fn-def} shows that the
common odd factor
\[
 (2\lfloor n/2\rfloor+1)!!
 \prod_{k=1}^{\lfloor n/2\rfloor}(2^{2k}-1)
\]
clears all their denominators.  This proves \eqref{eq:D-bound}.  The second
claim follows because the inverse-power value
$\upf(1-2^{-n})=F(2^{-n})$ is one of the values whose denominator divides
$D_n$.
\end{proof}

The following striking pattern is supported by computation and is recorded
as a conjecture, not as a theorem.

\begin{conjecture}[Denominator pattern]\label{conj:denominator}
Put $A_n=2^{-\binom n2}D_n$.  Then
\begin{enumerate}[label=\textup{(\alph*)}]
\item $A_{2n}=A_{2n+1}$ for $n\ge1$;
\item if $K_n=A_{2n-1}$, then $2(2n-1)!$ divides $K_n$;
\item $H_n=K_n/[2(2n-1)!]$ is an odd integer.
\end{enumerate}
\end{conjecture}

\subsection{Lucas parity and exact \texorpdfstring{$2$-adic}{2-adic} valuations}

Lucas's theorem modulo $2$ says that $\binom{N}{K}$ is odd exactly when every
binary digit of $K$ is bounded by the corresponding digit of $N$.  Since the
binary expansion of $2n+1$ is the expansion of $n$ followed by a final $1$,
we obtain:

\begin{proposition}\label{prop:odd-binomial-count}
For $n\ge1$, exactly $2^{\wt(n)}$ of the numbers
$\binom{2n+1}{2k}$, $0\le k\le n$, are odd.  Exactly
$2^{\wt(n)+1}$ of all the numbers $\binom{2n+1}{k}$ are odd.
\end{proposition}

\begin{theorem}[Oddness theorem]\label{thm:oddness}
For every $n\ge0$, $\mathcal F_n$ is odd.  For every $n\ge1$,
\begin{equation}\label{eq:2dn-valuation}
 \nu_2(2d_n)=0.
\end{equation}
Consequently every $R_n$ is odd, and
\begin{equation}\label{eq:up-valuation}
 \nu_2\bigl(\upf(1-2^{-n})\bigr)
 =-\binom n2-1-\nu_2(n!).
\end{equation}
\end{theorem}

\begin{proof}
Use \eqref{eq:Fn-recurrence} and induction.  All factors other than the
binomial coefficient are odd.  Among the $k=0,\ldots,n-1$ terms, the number
of odd binomial coefficients is $2^{\wt(n)}-1$, because the omitted
$k=n$ coefficient is odd.  This number is odd for $n\ge1$, so the sum is odd.

For odd indices, \eqref{eq:R-odd} and \eqref{eq:Rn-dn} imply
$\nu_2(2d_{2m+1})=0$.  For even indices, use the induction equation
\eqref{eq:2d-even-induction} after multiplying by a sufficiently large common
odd denominator.  By the induction hypothesis the terms with $k\ge1$ are
odd, while the $k=0$ term is even.  Proposition
\ref{prop:odd-binomial-count} implies that an odd number of the remaining
summands are odd, so the total is odd.  Hence \eqref{eq:2dn-valuation}.
Equation \eqref{eq:Rn-dn} then makes $R_n$ odd.  Finally take $2$-adic
valuations in \eqref{eq:d-F-inverse}.
\end{proof}

\subsection{Bernoulli recurrences}

Let $B_j$ be the Bernoulli numbers with $B_1=-1/2$.  The reciprocal
refinement factors
\[
 \frac{\pi z}{\sin\pi z},
 \qquad
 \frac{z}{\e^z-1}
\]
give alternative recurrences.

\begin{proposition}\label{prop:bernoulli-recurrences}
For $n\ge1$,
\begin{equation}\label{eq:c-bernoulli}
 c_n=\frac1{2^{2n}-1}
 \sum_{k=1}^{n}2^{2n-2k}(2^{2k}-2)
 \binom{2n}{2k}B_{2k}c_{n-k},
\end{equation}
and
\begin{equation}\label{eq:d-bernoulli}
 d_n=\frac{n2^{n-2}}{2^n-1}d_{n-1}
 -\frac1{2^n-1}
 \sum_{k=1}^{\lfloor n/2\rfloor}
 \binom{n}{2k}2^{n-2k}B_{2k}d_{n-2k}.
\end{equation}
\end{proposition}

\begin{proof}
Rewrite the Fourier and generating-function refinements as
\[
 \frac{\pi z}{\sin\pi z}\widehat{\upf}(z)=\widehat{\upf}(z/2),
 \qquad
 \frac{z}{\e^z-1}f(2z)=f(z).
\]
Insert the standard Bernoulli expansions of the two reciprocal factors and
compare coefficients.
\end{proof}

\section{The Fourier--Legendre expansion}\label{sec:legendre}

Because $\upf$ is even, its Legendre expansion contains only even degrees.
Let $P_m$ be the Legendre polynomial normalized by $P_m(1)=1$, and put
\begin{equation}\label{eq:legendre-coeff-def}
 u_n=\frac{4n+1}{2}\int_{-1}^{1}\upf(x)P_{2n}(x)\dd x.
\end{equation}

\begin{theorem}[Exact Legendre coefficients]\label{thm:legendre-coeff}
For every $n\ge0$,
\begin{equation}\label{eq:legendre-coeff}
 \begin{split}
 u_n={}&4^{-n}(4n+1)
 \sum_{k=0}^{n}(-1)^{n+k}
 \binom{2n}{n+k}\binom{2n+2k}{2n}\\
 &\hspace{35mm}\cdot(2k)!\,
 2^{\binom{2k+1}{2}}F(2^{-2k-1}).
 \end{split}
\end{equation}
Moreover,
\begin{equation}\label{eq:legendre-series}
 \upf(x)=\sum_{n=0}^{\infty}u_nP_{2n}(x)
\end{equation}
converges absolutely and uniformly on $[-1,1]$, including both endpoints.
\end{theorem}

\begin{proof}
The explicit monomial expansion
\begin{equation}\label{eq:P-even-monomial}
 P_{2n}(x)=4^{-n}\sum_{k=0}^{n}(-1)^{n+k}
 \binom{2n}{n+k}\binom{2n+2k}{2n}x^{2k}
\end{equation}
reduces \eqref{eq:legendre-coeff-def} to the even moments $c_k$.
From \eqref{eq:d-from-c} in odd degree, or directly from
\eqref{eq:d-integral} and evenness,
\[
 d_{2k+1}=\frac{2k+1}{2}c_k.
\]
Combining this with \eqref{eq:d-F-inverse} gives
\begin{equation}\label{eq:c-inverse-power}
 \frac{c_k}{2}=(2k)!\,2^{\binom{2k+1}{2}}F(2^{-2k-1}).
\end{equation}
Substitution into \eqref{eq:legendre-coeff-def} proves
\eqref{eq:legendre-coeff}.

For convergence, use the self-adjoint Legendre operator
\[
 \mathcal Lg=-\frac{\dd}{\dd x}\left((1-x^2)g'(x)\right),
 \qquad \mathcal LP_m=m(m+1)P_m.
\]
All derivatives of $\upf$ vanish at $\pm1$, so repeated integration by parts
has no boundary terms.  For every $r\ge1$,
\[
 \int_{-1}^{1}\upf P_m
 =\frac{1}{[m(m+1)]^r}
  \int_{-1}^{1}\mathcal L^r\upf\,P_m.
\]
Since $|P_m(x)|\le1$ on $[-1,1]$ and $\mathcal L^r\upf$ is continuous,
the ordinary Legendre coefficient is $O(m^{1-2r})$.  Taking $r\ge2$ makes
the coefficient series absolutely summable.  The Weierstrass test yields
absolute uniform convergence.  Orthogonality identifies its sum with
$\upf$ in $L^2$, and uniform convergence plus continuity upgrades the
identity to every point.
\end{proof}

\begin{remark}
Formula \eqref{eq:legendre-coeff} is an unusual bridge between a global
orthogonal-polynomial expansion and endpoint values at the sparse sequence
$2^{-1},2^{-3},2^{-5},\ldots$.  Every coefficient is rational, although the
series represents a transcendental-looking smooth bump.
\end{remark}

\part*{Endpoint asymptotics}
\addcontentsline{toc}{part}{Endpoint asymptotics}

\section{The negative Laplace product and dyadic renormalization}
\label{sec:asymptotics}

The endpoint $x\downarrow0$ is governed by the Laplace transform of the
random series $X$.  Put, for $s>0$,
\begin{equation}\label{eq:G-def}
 G(s)=f(-s)=\mathbb E\e^{-sX}.
\end{equation}
The refinement equation \eqref{eq:f-refinement} gives
\begin{equation}\label{eq:G-product}
 G(s)=\prod_{j=1}^{\infty}
 \frac{1-\e^{-s/2^j}}{s/2^j}.
\end{equation}
This is the basic product whose logarithm must be analyzed sharply.

\subsection{The exact Laplace identity}

\begin{proposition}\label{prop:laplace-F}
For $s>0$,
\begin{equation}\label{eq:laplace-F}
 \int_0^{\infty}\e^{-sx}F(x)\dd x=\frac{G(s)}{s}.
\end{equation}
Here $F(x)=1$ for $x\ge1$, as in the bounded convention.
\end{proposition}

\begin{proof}
The distribution function of $X$ is $F$.  Integration by parts on $[0,1]$
gives
\[
 \int_0^1\e^{-sx}F(x)\dd x
 =-\frac{\e^{-s}}s+\frac1s\int_0^1\e^{-sx}\dd F(x)
 =-\frac{\e^{-s}}s+\frac{G(s)}s.
\]
The constant tail contributes
$\int_1^\infty\e^{-sx}\dd x=\e^{-s}/s$, and the two boundary terms cancel.
\end{proof}

For every $r>0$, Bromwich inversion therefore has the form
\begin{equation}\label{eq:bromwich}
 F(x)=\frac{1}{2\pi\ii}
 \int_{r-\ii\infty}^{r+\ii\infty}
 \e^{sx}\frac{G(s)}s\dd s,
 \qquad x>0.
\end{equation}
The compact support of the underlying probability law and the rapid product
bounds justify the inversion and the contour manipulations used below.

\subsection{An exact periodic decomposition}

Define
\begin{equation}\label{eq:K-Laplace}
 K(u)=\log\left(\frac{1-\e^{-u}}u\right),
 \qquad
 \mathscr L(s)=\log G(s)=\sum_{j=1}^{\infty}K(s/2^j).
\end{equation}
The series is locally uniformly convergent for $s>0$.  It obeys
\begin{equation}\label{eq:L-refinement}
 \mathscr L(2s)-\mathscr L(s)=K(s).
\end{equation}
Let
\begin{equation}\label{eq:H-def}
 H(s)=\mathscr L(s)+\frac{(\log s)^2}{2L}-\frac{\log s}{2}.
\end{equation}
A direct calculation using $K(s)+\log s=\log(1-\e^{-s})$ gives
\begin{equation}\label{eq:H-defect}
 H(2s)-H(s)=\log(1-\e^{-s}).
\end{equation}
The defect is removed by the forward tail
\begin{equation}\label{eq:E-tail}
 E(s)=-\sum_{m=0}^{\infty}\log(1-\e^{-2^ms}).
\end{equation}
Indeed,
\begin{equation}\label{eq:E-defect}
 E(2s)-E(s)=\log(1-\e^{-s}).
\end{equation}
Consequently $H(s)-E(s)$ is invariant under $s\mapsto2s$.

\begin{theorem}[Exact dyadic decomposition]\label{thm:laplace-periodic}
There is a smooth one-periodic function $R:\R\to\R$ such that, for every
$s>0$,
\begin{equation}\label{eq:L-periodic-decomp}
 \mathscr L(s)
 =-\frac{(\log s)^2}{2L}+\frac{\log s}{2}
  +R(\log_2s)+E(s).
\end{equation}
The tail is positive and satisfies
\begin{equation}\label{eq:E-bound}
 0\le E(s)\le\frac{\e^{-s}}{(1-\e^{-s})^2}.
\end{equation}
In particular, $E(s)$ is exponentially small as $s\to\infty$ and is smaller
than every negative power of $\log s$.
\end{theorem}

\begin{proof}
Define $R(\log_2s)=H(s)-E(s)$.  Equations
\eqref{eq:H-defect} and \eqref{eq:E-defect} make this definition invariant
under replacing $s$ by $2s$, hence $R$ is one-periodic.  Smoothness follows
from locally uniform differentiation of the defining series.

For the bound, use $-\log(1-y)\le y/(1-y)$ and $2^m\ge m+1$:
\[
 E(s)\le\sum_{m\ge0}\frac{\e^{-2^ms}}{1-\e^{-2^ms}}
 \le\frac1{1-\e^{-s}}\sum_{m\ge0}\e^{-(m+1)s}
 =\frac{\e^{-s}}{(1-\e^{-s})^2}.
\]
Rearranging the definition of $R$ proves
\eqref{eq:L-periodic-decomp}.
\end{proof}

The term $R(\log_2s)$ is not an artifact of the proof.  It is the exact
record of the discrete dyadic scale in the infinite product.  Its mean and
all Fourier coefficients can be computed explicitly.

\section{The Gamma--zeta periodic fluctuation}\label{sec:periodic-fluctuation}

Set
\begin{equation}\label{eq:A0-def}
 A_0=\frac{6\gamma^2+12\gamma_1-\pi^2}{12}
 =\frac{\gamma^2}{2}+\gamma_1-\frac{\pi^2}{12}.
\end{equation}

\begin{theorem}[Mean and Fourier coefficients]\label{thm:psi-fourier}
The mean of the periodic function in \eqref{eq:L-periodic-decomp} is
\begin{equation}\label{eq:R-mean}
 \overline R=\int_0^1R(u)\dd u=\frac{A_0}{L}-\frac{L}{12}.
\end{equation}
Let
\begin{equation}\label{eq:Psi-def}
 \Psi(u)=R(u)-\overline R.
\end{equation}
Then $\Psi$ is real-valued, one-periodic, and has the absolutely convergent
Fourier series
\begin{equation}\label{eq:Psi-fourier}
 \Psi(u)=\sum_{k\in\Z\setminus\{0\}}\widehat\Psi(k)\e^{2\pi\ii ku},
\end{equation}
where
\begin{equation}\label{eq:Psi-coeff}
 \widehat\Psi(k)=
 -\frac1L\Gamma\left(-\frac{2\pi\ii k}{L}\right)
 \zeta\left(1-\frac{2\pi\ii k}{L}\right)
 \qquad(k\ne0).
\end{equation}
Every nonzero Fourier coefficient is nonzero.  In particular, $\Psi$ is not
constant and cannot be absorbed into an error term tending to zero.
\end{theorem}

\begin{proof}
The Mellin transform of $K$ in the fundamental strip $-1<\Re z<0$ is
\begin{equation}\label{eq:Mellin-K}
 \int_0^\infty K(u)u^{z-1}\dd u
 =-\Gamma(z)\zeta(1+z).
\end{equation}
One obtains this by integration by parts from
$K'(u)=(\e^u-1)^{-1}-u^{-1}$ and then by analytic continuation of the
classical Mellin integral for $(\e^u-1)^{-1}$.

Mellin inversion of the harmonic sum
$\sum_{j\ge1}K(s/2^j)$ introduces the geometric factor
$2^z/(1-2^z)$.  Moving the contour across the imaginary-axis poles
$z=2\pi\ii k/L$ yields the Fourier modes.  With the Fourier convention in
\eqref{eq:Psi-fourier}, the residue at the pole indexed by $-k$ is exactly
\eqref{eq:Psi-coeff}.

At $z=0$ there is a third-order pole.  The Laurent expansions
\[
 \Gamma(z)=\frac1z-\gamma+
 \left(\frac{\gamma^2}{2}+\frac{\pi^2}{12}\right)z+O(z^2),
\]
\[
 \zeta(1+z)=\frac1z+\gamma-\gamma_1z+O(z^2),
 \qquad
 \frac{2^z}{1-2^z}=-\frac1{Lz}-\frac12-\frac{Lz}{12}+O(z^3)
\]
give the quadratic and linear logarithmic terms already extracted in
\eqref{eq:L-periodic-decomp}; the remaining constant residue is
$A_0/L-L/12$, proving \eqref{eq:R-mean}.

Stirling's formula on vertical lines makes the Gamma factor exponentially
small in $|k|$, while $\zeta(1+\ii t)$ has at most subexponential growth.
Hence the Fourier series and all its termwise derivatives converge
absolutely.  The Gamma function has no zeros, and the classical zero-free
line $\zeta(1+\ii t)\ne0$ for $t\ne0$ shows that every coefficient in
\eqref{eq:Psi-coeff} is nonzero.
\end{proof}

Numerically,
\begin{align}
 \widehat\Psi(1)
 &\approx(7.5586467408-7.4700963646\,\ii)\times10^{-7},
 \label{eq:first-psi-numeric}\\
 |\widehat\Psi(1)|&\approx1.0627110626\times10^{-6}.
\end{align}
The second harmonic has magnitude about $6.69\times10^{-13}$ and subsequent
ones are dramatically smaller.  Thus the oscillation is visually tiny, but
its amplitude does not decay as $x\downarrow0$; asymptotically, that is what
matters.

\section{The exact lower-Lambert phase}\label{sec:lambert-phase}

The saddle scale is not simply $\log_2(1/x)$.  It is the large solution of a
Lambert equation.

\begin{definition}\label{def:lambda}
For $0<x<1/(\e L)$, define
\begin{equation}\label{eq:lambda-W}
 \lambda(x)=-\frac1L W_{-1}(-Lx),
\end{equation}
where $W_{-1}$ is the lower real branch of Lambert's $W$-function.
\end{definition}

The defining identity $W\e^W=z$ gives
\begin{equation}\label{eq:lambda-equation}
 \lambda(x)2^{-\lambda(x)}=x.
\end{equation}
Thus, with
\begin{equation}\label{eq:radius}
 r(x)=2^{\lambda(x)},
\end{equation}
we have the exact saddle relation
\begin{equation}\label{eq:rx-lambda}
 r(x)x=\lambda(x),
 \qquad
 r(x)=\frac{\lambda(x)}x.
\end{equation}
The lower branch is essential: it selects the solution
$\lambda>1/L$ and hence $\lambda(x)\to\infty$ as $x\downarrow0$.

If $x=2^{-t}$, then \eqref{eq:lambda-equation} is equivalent to the fixed
point equation
\begin{equation}\label{eq:lambda-fixed-point}
 \lambda-\frac{\log\lambda}{L}=t.
\end{equation}
In particular,
\begin{equation}\label{eq:lambda-sim-t}
 \lambda(t)\sim t,
 \qquad
 \lambda(t)-t\sim\frac{\log t}{L}.
\end{equation}
This distinction is large enough to affect several lower-order terms in
$\log F(2^{-t})$.

\section{The sharp endpoint main term}\label{sec:sharp-main}

Define the constant
\begin{equation}\label{eq:Csharp}
 C_*=\frac{A_0}{L}-\frac{7L}{12}-\frac12\log\pi
 =\frac{6\gamma^2+12\gamma_1-\pi^2}{12L}
  -\frac{7L}{12}-\frac12\log\pi.
\end{equation}
For sufficiently small $x>0$, set
\begin{equation}\label{eq:Msharp}
 \begin{split}
 M(x)={}&-\frac L2\lambda(x)^2
 +\left(1+\frac L2\right)\lambda(x)
 -\frac12\log\lambda(x)\\
 &\quad+C_*+\Psi(\lambda(x)).
 \end{split}
\end{equation}

\begin{theorem}[Sharp endpoint asymptotic]\label{thm:sharp-main}
As $x\downarrow0$,
\begin{equation}\label{eq:sharp-log}
 \log F(x)=M(x)+O\left(\frac1{\lambda(x)}\right)
          =M(x)+O\left(\frac1{-\log x}\right).
\end{equation}
More precisely,
\begin{equation}\label{eq:sharp-ratio}
 F(x)=\e^{M(x)}\left(1+O\left(\frac1{\lambda(x)}\right)\right).
\end{equation}
The periodic function in \eqref{eq:Msharp} is evaluated at the exact phase
\eqref{eq:lambda-W}.
\end{theorem}

\begin{proof}
We describe the saddle argument in enough detail to locate every term in
\eqref{eq:Msharp}.  Start from Bromwich inversion \eqref{eq:bromwich}, choose
$r=2^\lambda$ with $rx=\lambda$, and parametrize the vertical line by
\[
 s=r(1+\ii y),\qquad y\in\R.
\]
Since $\dd s/s=\ii\dd y/(1+\ii y)$,
\begin{equation}\label{eq:bromwich-y}
 F(x)=\frac1{2\pi}\int_{-\infty}^{\infty}
 \frac{\exp\{\lambda(1+\ii y)+\mathscr L(r(1+\ii y))\}}
      {1+\ii y}\dd y.
\end{equation}
The real decomposition \eqref{eq:L-periodic-decomp} extends analytically to
the part of the right half-plane used by the central contour.  Put
$\delta(y)=\Log(1+\ii y)$.  Then
\[
 \Log s=L\lambda+\delta(y),
 \qquad
 \log_2s=\lambda+\frac{\delta(y)}L.
\]
Ignoring for the moment the exponentially small forward tail, the logarithm
of the integrand, including the factor $(1+\ii y)^{-1}$, is
\begin{equation}\label{eq:central-exponent}
 \begin{split}
 \mathcal E_\lambda(y)={}&\lambda(1+\ii y)
 -\frac{(L\lambda+\delta(y))^2}{2L}
 +\frac{L\lambda+\delta(y)}2\\
 &+\overline R
 +\Psi\left(\lambda+\frac{\delta(y)}L\right)-\delta(y).
 \end{split}
\end{equation}
At $y=0$ this equals
\begin{equation}\label{eq:central-height}
 -\frac L2\lambda^2+
 \left(1+\frac L2\right)\lambda+
 \overline R+\Psi(\lambda).
\end{equation}
The choice $rx=\lambda$ cancels the linear term in $y$.  Since
\[
 \delta(y)=\ii y+\frac{y^2}{2}-\frac{\ii y^3}{3}
             -\frac{y^4}{4}+O(y^5),
\]
the real part of the exponent begins with
$-\lambda y^2/2$.  On the central scale $y=v/\sqrt\lambda$, the normalized
integral therefore tends to
\[
 \frac1{2\pi\sqrt\lambda}
 \int_{-\infty}^{\infty}\e^{-v^2/2}\dd v
 =\frac1{\sqrt{2\pi\lambda}}.
\]
Taking logarithms contributes
$-\frac12\log\lambda-\frac12\log(2\pi)$.  Since
\[
 \overline R-\frac12\log(2\pi)
 =\frac{A_0}{L}-\frac{L}{12}-\frac L2-\frac12\log\pi=C_*,
\]
we recover exactly \eqref{eq:Msharp}.

For completeness, the error analysis has three parts.  On a shrinking central
arc, Taylor expansion of \eqref{eq:central-exponent}, dominated by a fixed
Gaussian, gives a relative error $O(1/\lambda)$.  On the adjacent arc, the
real part is bounded above by $-c\lambda y^2$.  On the remaining minor arc,
blocks of factors in the product \eqref{eq:G-product} give uniform geometric
suppression; the resulting integral is exponentially smaller than the
central mass.  Finally, \eqref{eq:E-bound} at $r=2^\lambda$ shows that the
forward tail is exponentially flat in $\lambda$ and does not affect any
inverse-power term.  These estimates justify contour truncation, termwise
expansion, and the passage from the relative estimate
\eqref{eq:sharp-ratio} to \eqref{eq:sharp-log}.
\end{proof}

\subsection{The leading quadratic law}

\begin{corollary}\label{cor:quadratic-law}
As $x\downarrow0$,
\begin{equation}\label{eq:quadratic-law-x}
 \frac{\log F(x)}{(\log x)^2}\longrightarrow-\frac1{2\log2}.
\end{equation}
Equivalently, as $t\to\infty$,
\begin{equation}\label{eq:quadratic-law-t}
 \frac{\log F(2^{-t})}{t^2}\longrightarrow-\frac{\log2}{2}.
\end{equation}
Consequently, for every $A,\alpha>0$,
\begin{equation}\label{eq:flat-comparison}
 F(x)=o(x^A),
 \qquad
 \e^{-x^{-\alpha}}=o(F(x)).
\end{equation}
\end{corollary}

\begin{proof}
The phase satisfies $\lambda\sim-\log x/L$, and every other term in
\eqref{eq:Msharp} is $o(\lambda^2)$.  This gives
\eqref{eq:quadratic-law-x} and \eqref{eq:quadratic-law-t}.  The comparisons in
\eqref{eq:flat-comparison} follow by comparing the quadratic logarithmic
scale with $\log x$ and with $-x^{-\alpha}$, respectively.
\end{proof}

\section{Corrected explicit forms and the failure of the nonperiodic formula}
\label{sec:corrected-explicit}

\subsection{The compact Lambert--\texorpdfstring{$W$}{W} form}

Let $W=W_{-1}(-Lx)$.  Since $W=-L\lambda(x)$ and
$\log x=\log\lambda-L\lambda$, the main term
\eqref{eq:Msharp} is exactly
\begin{equation}\label{eq:compact-W-main}
 \boxed{
 M(x)=-\frac{7L}{12}-\frac12\log(\pi x)
 +\frac{A_0-W-W^2/2}{L}
 +\Psi\left(-\frac WL\right).}
\end{equation}
This is the cleanest closed expression for the corrected leading asymptotic.
Theorem \ref{thm:sharp-main} says that
\begin{equation}\label{eq:compact-W-asymptotic}
 \log F(x)=\text{right side of \eqref{eq:compact-W-main}}
 +O\left(\frac1{-\log x}\right).
\end{equation}

\subsection{An elementary logarithmic expansion}

To compare with formulas written without Lambert $W$, set
\begin{equation}\label{eq:ellma}
 \ell=\log x<0,
 \qquad m=\log(-\ell),
 \qquad a=\log L.
\end{equation}
Define
\begin{equation}\label{eq:E0-explicit}
 \begin{split}
 \mathcal E_0(x)={}&-\frac{\ell^2}{2L}+\frac{\ell m}{L}
 -\left(\frac12+\frac{1+a}{L}\right)\ell
 -\frac{m^2}{2L}+\frac{am}{L}\\
 &+\frac{6\gamma^2+12\gamma_1-\pi^2-6a^2}{12L}
 -\frac{7L}{12}-\frac12\log\pi.
 \end{split}
\end{equation}
Expansion of the exact phase gives
\begin{equation}\label{eq:elementary-corrected-coarse}
 \log F(x)=\mathcal E_0(x)+\Psi(\lambda(x))
 +O\left(\frac{m^2}{-\ell}\right).
\end{equation}
Keeping the first inverse-logarithmic phase term gives
\begin{equation}\label{eq:E1-explicit}
 \mathcal E_1(x)=\mathcal E_0(x)+\frac{(m-a)^2}{2L\ell}
\end{equation}
and
\begin{equation}\label{eq:elementary-corrected-sharp}
 \log F(x)=\mathcal E_1(x)+\Psi(\lambda(x))
 +O\left(\frac1{-\ell}\right).
\end{equation}
The sign in the last term is encoded by $\ell<0$.

\subsection{Dyadic logarithmic scale}

For $x=2^{-t}$, write $u=\log t$.  Then the exact phase satisfies
\eqref{eq:lambda-fixed-point}, and the main term expands as
\begin{equation}\label{eq:dyadic-main-expanded}
 \begin{split}
 M(2^{-t})={}&-\frac L2t^2-tu+
 \left(1+\frac L2\right)t
 -\frac{u^2}{2L}+C_*+\Psi(\lambda(t))\\
 &-\frac{u^2}{2L^2t}
 +O\left(\frac{u^3}{t^2}\right).
 \end{split}
\end{equation}
Combining this with \Cref{thm:sharp-main} gives a convenient corrected
formula with error $O(1/t)$.

\subsection{Why the periodic term cannot be omitted}

Let
\begin{equation}\label{eq:nonperiodic-main}
 M_0(x)=M(x)-\Psi(\lambda(x)).
\end{equation}
Then \Cref{thm:sharp-main} gives
\begin{equation}\label{eq:omission-error}
 \log F(x)-M_0(x)=\Psi(\lambda(x))+o(1).
\end{equation}
Choose any $u\in[0,1)$ and the sequence
\begin{equation}\label{eq:phase-subsequence}
 \lambda_N=N+u,
 \qquad x_N=\lambda_N2^{-\lambda_N}.
\end{equation}
Then $x_N\downarrow0$ and the exact phase of $x_N$ is $\lambda_N$, so
\begin{equation}\label{eq:phase-limit}
 \log F(x_N)-M_0(x_N)\longrightarrow\Psi(u).
\end{equation}
Because $\Psi$ is nonconstant, these subsequential limits are not all equal.
Therefore:
\begin{enumerate}
\item the nonperiodic logarithmic main term does not have an $o(1)$ error;
\item the quotient $F(x)/\e^{M_0(x)}$ does not tend to $1$;
\item no claimed error such as $O((\log x)^{-1})$ can hold after simply
      deleting $\Psi$.
\end{enumerate}
The first harmonic is only about one part in a million, but a fixed one-part-
in-a-million oscillation is still much larger than every error tending to
zero.

\section{The first saddle correction}\label{sec:first-correction}

The $O(1/\lambda)$ in \Cref{thm:sharp-main} has an explicit periodic
coefficient.  Let primes denote ordinary derivatives of the one-periodic
function $\Psi$.  Define
\begin{equation}\label{eq:A-B-saddle}
 A(u)=\frac{\Psi'(u)}L-\frac12,
 \qquad
 B(u)=-\frac14+\frac1{2L}+\frac{\Psi'(u)}{2L}
      -\frac{\Psi''(u)}{2L^2}.
\end{equation}
On the central scale $y=v/\sqrt\lambda$, the first exponent polynomials are
\begin{equation}\label{eq:P1P2}
 P_1(v;u)=\ii\left(A(u)v+\frac{v^3}{3}\right),
 \qquad
 P_2(v;u)=\frac{v^4}{4}+B(u)v^2.
\end{equation}
They come directly from expanding \eqref{eq:central-exponent}; the cubic term
is a characteristic feature of the vertical Bromwich parametrization
$s=r(1+\ii y)$.

\begin{definition}\label{def:c1}
The first logarithmic saddle coefficient is
\begin{equation}\label{eq:c1}
 c_1(u)=\frac1{24}+\frac1{2L}
 -\frac{\Psi''(u)+\Psi'(u)^2}{2L^2}.
\end{equation}
\end{definition}
It is continuous, bounded, smooth, and one-periodic.

\begin{theorem}[First correction]\label{thm:first-correction}
As $x\downarrow0$,
\begin{equation}\label{eq:first-correction-log}
 \log F(x)=M(x)+\frac{c_1(\lambda(x))}{\lambda(x)}
 +O\left(\frac1{\lambda(x)^2}\right).
\end{equation}
Equivalently,
\begin{equation}\label{eq:first-correction-ratio}
 F(x)=\e^{M(x)}\left[
 1+\frac{c_1(\lambda(x))}{\lambda(x)}
 +O\left(\frac1{\lambda(x)^2}\right)
 \right].
\end{equation}
\end{theorem}

\begin{proof}
After factoring out the height \eqref{eq:central-height}, setting
$y=v/\sqrt\lambda$, and retaining terms through $\lambda^{-1}$, the central
integrand has the form
\begin{equation}\label{eq:central-first-expansion}
 \e^{-v^2/2}\left[
 1+\frac{P_1(v;u)}{\sqrt\lambda}
 +\frac{P_2(v;u)+\frac12P_1(v;u)^2}{\lambda}
 +O\left(\frac{(1+|v|)^9}{\lambda^{3/2}}
\right)
 \right],
\end{equation}
where $u=\lambda$ modulo $1$.  The odd term integrates to zero.  If $V$ is a
standard normal random variable, the relative coefficient of $1/\lambda$ is
\begin{equation}\label{eq:gaussian-contraction}
 \mathbb E\left[P_2(V;u)+\frac12P_1(V;u)^2\right].
\end{equation}
Using
$\mathbb EV^2=1$, $\mathbb EV^4=3$, and $\mathbb EV^6=15$, followed by the
substitutions \eqref{eq:A-B-saddle}, reduces
\eqref{eq:gaussian-contraction} to
\[
 \frac1{24}+\frac1{2L}
 -\frac{\Psi''(u)+\Psi'(u)^2}{2L^2}=c_1(u).
\]
The minor arcs and the forward tail are smaller than $\lambda^{-2}$ relative
to the main mass after a one-order refinement of the estimates in
\Cref{thm:sharp-main}.  This proves \eqref{eq:first-correction-ratio}; taking
the logarithm leaves the same first coefficient and proves
\eqref{eq:first-correction-log}.
\end{proof}

Combining \eqref{eq:dyadic-main-expanded} and
\eqref{eq:first-correction-log} gives the particularly useful formula
\begin{equation}\label{eq:dyadic-first-full}
 \begin{split}
 \log F(2^{-t})={}&-\frac L2t^2-t\log t
 +\left(1+\frac L2\right)t
 -\frac{(\log t)^2}{2L}+C_*+\Psi(\lambda(t))\\
 &-\frac{(\log t)^2}{2L^2t}
 +\frac{c_1(\lambda(t))}{\lambda(t)}
 +O\left(\frac{(\log t)^3}{t^2}\right).
 \end{split}
\end{equation}
The periodic functions remain at the exact phase; only the elementary part
has been expanded in $t$.

\section{The full all-orders expansion}\label{sec:full-expansion}

The saddle method does not stop at the first correction.  The periodic
Fourier series is differentiable to all orders, the forward tail is
exponentially flat, and the central kernel admits a complete Gaussian
expansion.

\begin{theorem}[Full exact-phase expansion]\label{thm:full-expansion}
There exist bounded smooth one-periodic functions
$c_j:\R\to\R$ for $j\ge0$, with
\begin{equation}\label{eq:c0-zero}
 c_0=0,
 \qquad c_1\text{ given by \eqref{eq:c1}},
\end{equation}
such that, for every integer $N\ge1$,
\begin{equation}\label{eq:full-expansion}
 \log F(x)=M(x)+
 \sum_{j=1}^{N-1}\frac{c_j(\lambda(x))}{\lambda(x)^j}
 +O\left(\frac1{\lambda(x)^N}\right)
 \qquad(x\downarrow0).
\end{equation}
Equivalently, the error is
$O(({-\log x})^{-N})$.  The coefficient functions are evaluated at the exact
lower-Lambert phase.
\end{theorem}

\begin{proof}
We outline the all-orders construction.

\emph{Central analytic expansion.}
In \eqref{eq:central-exponent}, put $y=v/\sqrt\lambda$ and expand
$\delta(y)=\Log(1+\ii y)$, the analytic periodic function
$\Psi(\lambda+\delta/L)$, and the amplitude $(1+\ii y)^{-1}$ to an arbitrary
fixed order.  After the Gaussian term $-v^2/2$ is removed, one obtains
exponent polynomials
\begin{equation}\label{eq:exponent-polynomials}
 \sum_{m=1}^{2N}\lambda^{-m/2}P_m(v;\lambda),
\end{equation}
where every coefficient of $P_m$ is a polynomial in finitely many derivatives
of $\Psi$ evaluated at $\lambda$.  These coefficients are bounded and
one-periodic.

\emph{Exponentiation and Gaussian contraction.}
Expanding the exponential of \eqref{eq:exponent-polynomials} produces finite
Bell-polynomial combinations.  Integrating termwise against
$\e^{-v^2/2}$ kills every odd total power.  Thus only integral powers
$\lambda^{-j}$ survive.  The resulting relative-mass coefficients
$b_j(\lambda)$ are bounded smooth periodic functions, with $b_0=1$.

\emph{Uniform remainder.}
Choose a central window that grows as a small power of $\lambda$ in the
$v$-variable.  Taylor's theorem and the exponential decay of the Gamma--zeta
Fourier coefficients give an integrable Gaussian majorant for the truncated
remainder.  The intermediate and minor arcs satisfy estimates stronger than
any prescribed inverse power after the truncation parameters are chosen in
the usual order.  The tail $E(2^\lambda)$ is exponentially small and may be
added to an expansion of every finite order.

\emph{Taking the logarithm.}
The saddle ratio has an expansion
$1+\sum_{j=1}^{N-1}b_j(\lambda)\lambda^{-j}+O(\lambda^{-N})$.
The formal power series for $\log(1+z)$ converts the $b_j$ into periodic
polynomials $c_j$ without changing the remainder order.  The first one is
$c_1=b_1$, already computed in \Cref{thm:first-correction}.  This proves
\eqref{eq:full-expansion}.
\end{proof}

\subsection{All-order expansion of the phase}

The exact phase can itself be expanded in powers of $1/t$.  Let $u=\log t$.
Solving \eqref{eq:lambda-fixed-point} recursively gives
\begin{equation}\label{eq:phase-first-terms}
 \begin{split}
 \lambda(t)={}&t+\frac{u}{L}
 +\frac{u}{L^2t}
 +\frac{u-u^2/2}{L^3t^2}\\
 &+\frac{u-3u^2/2+u^3/3}{L^4t^3}
 +O\left(\frac{(1+u)^4}{t^4}\right).
 \end{split}
\end{equation}
More generally, there are unique polynomials $P_j(u)$ such that, for every
$N$,
\begin{equation}\label{eq:phase-all-orders}
 \lambda(t)=t+\frac{u}{L}+
 \sum_{j=1}^{N-1}\frac{P_j(u)}{t^j}
 +O\left(\frac{(1+u)^N}{t^N}\right).
\end{equation}
They are obtained by substituting the right side into
$\lambda=t+L^{-1}\log\lambda$ and cancelling successive powers of $1/t$.

\begin{remark}[Exact phase versus expanded phase]
Theorem \ref{thm:full-expansion} deliberately keeps every periodic
coefficient at $\lambda(x)$, not at a truncated approximation to it.  One may
combine \eqref{eq:phase-all-orders} with Taylor expansions of $\Psi$ and the
$c_j$ to any fixed order.  A single all-orders statement after composing all
periodic functions with a formal phase expansion requires a separate
composition theorem and additional bookkeeping; it is not silently assumed
here.
\end{remark}

\subsection{A numerical check at exact inverse powers}

The values $F(2^{-n})$ are exact rationals, so they provide unusually clean
tests of the asymptotic formula.  The following table compares the exact
logarithm with $M$; the last column is the predicted first coefficient at the
corresponding exact phase.
\begin{center}
\begin{tabular}{@{}rrrrr@{}}
\toprule
$n$ & $\lambda(2^{-n})$ & $\log F(2^{-n})-M$ &
$\lambda(\log F-M)$ & $c_1(\lambda)$\\
\midrule
$10$ & $13.78503056$ & $0.0580133344$ & $0.79971559$ & $0.76296785$\\
$20$ & $24.62186833$ & $0.0318275697$ & $0.78365423$ & $0.76292687$\\
$40$ & $45.50804986$ & $0.0170138895$ & $0.77426893$ & $0.76294905$\\
\bottomrule
\end{tabular}
\end{center}
The scaled residual moves toward the periodic first correction; the remaining
difference is of order $1/\lambda$, as predicted by
\eqref{eq:first-correction-log}.

\section{Consequences and conceptual summary}\label{sec:consequences}

Several themes now fit together.

\paragraph{Self-similarity versus exact arithmetic.}
The global equation $\Th'(x)=2\Th(2x)$ turns differentiation into a binary
magnification.  At dyadic points, sufficient magnification lands at an
integer, where $\Th$ is either zero or a Thue--Morse sign.  This is why Taylor
data terminate and why dyadic values are rational.

\paragraph{The moments are the bridge.}
The same $c_n$ appear as Fourier coefficients near the origin, complete even
moments of the bump, coefficients in the exact dyadic formula, and building
blocks of the Legendre expansion.  Their integer normalization
$\mathcal F_n$ controls denominators and parity.

\paragraph{The endpoint remembers the discrete scale.}
The quadratic term in $\log F(x)$ comes from the triple pole at zero in the
Mellin representation of the dyadic harmonic sum.  The other poles on the
imaginary axis encode the one-periodic fluctuation.  The fluctuation is tiny
because Gamma decays exponentially on vertical lines, but it is nonzero
because neither Gamma nor zeta vanishes at the relevant points.

\paragraph{Lambert $W$ is not cosmetic.}
The equation $\lambda2^{-\lambda}=x$ exactly balances the Laplace radius and
the endpoint argument.  Replacing $\lambda$ too early by
$\log_2(1/x)$ loses terms involving $\log\log(1/x)$ and misplaces the phase of
$\Psi$.  The exact lower branch is therefore the natural coordinate for the
full expansion.

\paragraph{Smoothness is extreme but intermediate.}
The Fabius function is flatter than every power at the endpoint but much
larger than $\exp(-x^{-\alpha})$ for every $\alpha>0$.  Its Taylor series is
zero at the endpoint and finite at every reduced dyadic point, yet it is
nowhere analytic.  The asymptotic formula supplies a quantitative profile for
this otherwise qualitative pathology.


\appendix

\section{Infinite binary-reduction series and discrete limits}
\label{app:global-series}

The finite dyadic formulas of the main text terminate because the binary tail
of a dyadic number eventually vanishes.  For a general real argument the same
Taylor reductions do not terminate, but they telescope.  This appendix states
the resulting absolutely convergent series, including the scale-zero term
that is easy to overlook, and then records the related discrete-limit
formula.

\subsection{Binary prefixes, tails, and the analytic telescope}

For $x\geq0$ and $m\in\N$, put
\begin{equation}\label{eq:binary-prefix-tail}
 A_m=A_m(x)=\lfloor 2^m x\rfloor,
 \qquad
 Y_m=Y_m(x)=x-\frac{A_m}{2^m},
\end{equation}
and define the previous prefix by
\begin{equation}\label{eq:binary-previous-prefix}
 B_0=B_0(x)=\left\lfloor\frac{x}{2}\right\rfloor,
 \qquad
 B_m=B_m(x)=A_{m-1}(x)\quad(m\geq1).
\end{equation}
Thus $0\leq Y_m<2^{-m}$.  The integer
$A_m-2B_m\in\{0,1\}$ is the $m$th newly exposed binary digit, with the
scale-zero convention built into \eqref{eq:binary-previous-prefix}.

Define the finite reduction polynomial
\begin{equation}\label{eq:reduction-polynomial}
 \mathcal R_m(y)=
 \sum_{k=0}^{m}
 2^{\binom{k+1}{2}-\binom{m-k}{2}}
 \frac{d_{m-k}}{(m-k)!}\frac{y^k}{k!}.
\end{equation}
This is exactly the Taylor polynomial appearing in
\eqref{eq:taylor-recursion}, now indexed by the scale $m$.

\begin{theorem}[Global binary-reduction series]
\label{thm:global-binary-series}
For every $x\geq0$,
\begin{equation}\label{eq:global-binary-series}
 \Th(x)=\sum_{m=0}^{\infty}
 \epsilon_{A_m}\,(2B_m-A_m)\,\mathcal R_m(Y_m),
\end{equation}
and the series is absolutely convergent.  If $0\leq x\leq1$, the left side
may be replaced by the bounded Fabius function $F(x)$.
\end{theorem}

\begin{proof}
Let
\begin{equation}\label{eq:binary-remainder}
 E_m(x)=\epsilon_{A_m}\Th(Y_m).
\end{equation}
The one-step Taylor reduction says that passing from the $(m-1)$st tail to
the $m$th tail either changes nothing, when the new binary digit is $0$, or
crosses one inverse-power center and contributes the polynomial
$\mathcal R_m$.  In both cases the identity can be written uniformly as
\begin{equation}\label{eq:binary-one-step}
 E_{m-1}(x)-E_m(x)
 =\epsilon_{A_m}(2B_m-A_m)\mathcal R_m(Y_m)
 \qquad(m\geq1).
\end{equation}
At scale zero, the same decomposition is valid after first locating $x$ in
its even or odd block of length one:
\begin{equation}\label{eq:binary-zero-step}
 \Th(x)=\epsilon_{A_0}(2B_0-A_0)\mathcal R_0(Y_0)+E_0(x).
\end{equation}
Summing \eqref{eq:binary-one-step} and using
\eqref{eq:binary-zero-step} gives the exact finite telescope
\begin{equation}\label{eq:binary-finite-telescope}
 \Th(x)=\sum_{m=0}^{N}
 \epsilon_{A_m}(2B_m-A_m)\mathcal R_m(Y_m)+E_N(x).
\end{equation}
Since $0\leq Y_N<2^{-N}$ and $F(y)\leq2y$ on $[0,1/2]$, for $N\geq1$ we
have
\begin{equation}\label{eq:binary-remainder-bound}
 |E_N(x)|\leq2^{1-N}.
\end{equation}
Hence $E_N(x)\to0$, proving ordinary convergence.  Moreover,
\eqref{eq:binary-one-step} and the same bound give, for $m\geq2$,
\begin{equation}\label{eq:binary-summand-bound}
 \bigl|\epsilon_{A_m}(2B_m-A_m)\mathcal R_m(Y_m)\bigr|
 \leq4\,2^{-(m-1)}.
\end{equation}
The two initial terms are finite and the geometric majorant proves absolute
convergence.  On the closed unit interval, $\Th=F$.
\end{proof}

\subsection{The parity-power form and the missing endpoint term}

The inverse-power bridge \eqref{eq:d-F-inverse} rewrites half of the reduction
polynomial as
\begin{equation}\label{eq:parity-inner}
 \mathcal I_m(y)=
 \sum_{n=0}^{m}
 2^{\binom{n+1}{2}-1}\,
 \Th\bigl(2^{n-m}\bigr)\frac{y^n}{n!}
 =\frac12\mathcal R_m(y).
\end{equation}
The exponent for $n=0$ is genuinely $-1$.  With the usual convention
$0^0=1$, no exceptional interpretation is required.

\begin{corollary}[Parity-power series]
\label{cor:parity-power-series}
For $x\geq0$,
\begin{equation}\label{eq:parity-power-series}
 \Th(x)=\sum_{m=0}^{\infty}
 \bigl((-1)^{A_m}-1\bigr)\epsilon_{A_m}\,\mathcal I_m(Y_m),
\end{equation}
and the series is absolutely convergent.
\end{corollary}

\begin{proof}
For every nonnegative integer $A$,
\begin{equation}\label{eq:parity-coefficient-identity}
 \frac{(-1)^A-1}{2}=2\left\lfloor\frac A2\right\rfloor-A.
\end{equation}
Since $B_m=\lfloor A_m/2\rfloor$, including $m=0$, multiplication of
\eqref{eq:parity-inner} by \eqref{eq:parity-coefficient-identity} turns the
summand in \eqref{eq:parity-power-series} into the summand in
\eqref{eq:global-binary-series}.
\end{proof}

\begin{remark}[Why the index must begin at zero]
For $0\leq x<1$, $A_0=0$, so the $m=0$ term in
\eqref{eq:parity-power-series} vanishes.  This explains why a formerly
one-indexed version happens to work on the half-open interval.  At $x=1$,
however,
\[
 A_0=1,\qquad Y_0=0,\qquad
 \mathcal I_0(0)=2^{-1}\Th(1)=\frac12,
\]
and the omitted term equals
$((-1)^1-1)\epsilon_1\mathcal I_0(0)=1$.  Thus the scale-zero term is
indispensable at the right endpoint and, globally, distinguishes the even and
odd unit blocks of $\Th$.
\end{remark}

\subsection{The fully explicit global \texorpdfstring{$q$}{q}-binomial series}

For a scalar $q\in\R$ or $q\in\C$, define
\begin{equation}\label{eq:q-global-Cn}
 \mathcal C_n(q)=
 \frac{1}{2^{\binom n2}(1/2;1/2)_n}
 \sum_{k=0}^{n}
 \frac{\qbinom{n}{k}}{4^{\binom k2}(n+k)!}
 \sum_{r=0}^{2^k-1}
 \epsilon_r\,(r-2^k+q)^{n+k},
\end{equation}
and
\begin{equation}\label{eq:q-global-polynomial}
 \mathcal P_m(q;z)=
 \frac{1}{2^{\binom{m+1}{2}}}
 \sum_{n=0}^{m}\frac{z^{m-n}}{(m-n)!}\mathcal C_n(q).
\end{equation}
Although the right side is written using $q$, it is a constant polynomial in
that common translation.

\begin{theorem}[Global $q$-binomial--Thue--Morse series]
\label{thm:global-q-series}
Let $x\geq0$, and let $q$ be any real or complex number.  Then
\begin{equation}\label{eq:global-q-series}
 \Th(x)=\sum_{m=0}^{\infty}
 \epsilon_{A_m}(2B_m-A_m)
 \mathcal P_m\bigl(q;2^{m+1}Y_m\bigr).
\end{equation}
For complex $q$ the equality is understood after embedding the real number
$\Th(x)$ into $\C$.  The series is absolutely convergent, and every single
summand is independent of $q$ before summation.  On $[0,1]$ the sum equals
$F(x)$.
\end{theorem}

\begin{proof}
The finite translation-invariance argument of
\Cref{sec:qbinomial} applies to the inner polynomial at each fixed scale and
gives the pointwise identity
\begin{equation}\label{eq:q-global-pointwise}
 \mathcal P_m\bigl(q;2^{m+1}y\bigr)=\mathcal R_m(y).
\end{equation}
One way to see the normalization is to expand the left side by
\eqref{eq:q-global-Cn}--\eqref{eq:q-global-polynomial}: the Gaussian-binomial
annihilator removes every translation term of degree below the active order,
and the surviving coefficient is precisely the inverse-power Taylor datum in
\eqref{eq:reduction-polynomial}.  Substitution of
\eqref{eq:q-global-pointwise} into \eqref{eq:global-binary-series} proves the
identity and imports absolute convergence.  Since the equality is polynomial,
the rational proof extends without change to real and complex translations.
\end{proof}

\subsection{A finite discrete approximation converging to \texorpdfstring{$\Th$}{Theta}}

For $x\geq0$, let
\begin{equation}\label{eq:discrete-range}
 N_p(x)=\left\lfloor 2^p x+\frac12\right\rfloor.
\end{equation}
Equivalently, the inner sum below has inclusive upper endpoint
$\lfloor2^p x-1/2\rfloor$; the formulation using $N_p$ also covers the empty
case safely.  For $q\in\R$ or $q\in\C$, define
\begin{align}
 D_n(q,x)
 &=\frac{1}{2^{n^2}(1/2;1/2)_n}
 \sum_{k=0}^{n}
 \frac{\genfrac{[}{]}{0pt}{}{n}{k}_{1/2}}{4^{\binom k2}(n+k)!}
 \notag\\[-2pt]
 &\hspace{22mm}\times
 \sum_{r=0}^{N_{n+k}(x)-1}
 \epsilon_r\bigl(r-2^{n+k}x+q\bigr)^{n+k}.
 \label{eq:discrete-limit-approximation}
\end{align}
An empty inner range contributes zero.

\begin{theorem}[Discrete-limit formula]
\label{thm:discrete-limit}
For every fixed $x\geq0$ and every fixed real or complex $q$,
\begin{equation}\label{eq:discrete-limit}
 \lim_{n\to\infty}D_n(q,x)=\Th(x).
\end{equation}
In particular, the limit is $F(x)$ for $0\leq x\leq1$ and is independent of
the translation $q$.
\end{theorem}

\begin{proof}
The proof separates an exact algebraic reindexing from two convergence
estimates.

Put $j=n-k$.  After normalizing the finite Thue--Morse power sum as a spline
branch of degree $2n-j$, the outer coefficients become the triangular array
\begin{equation}\label{eq:toeplitz-weight}
 w_{n,j}=
 \frac{(-1)^j\qbinom{n}{j}
       2^{-\binom{j+1}{2}}}{(1/2;1/2)_n},
 \qquad 0\leq j\leq n.
\end{equation}
The finite $q$-binomial theorem gives the exact row-sum identity
\begin{equation}\label{eq:toeplitz-row-sum}
 \sum_{j=0}^{n}w_{n,j}=1.
\end{equation}
The same theorem with the absolute signs removed only after taking norms
shows the uniform variation bound
\begin{equation}\label{eq:toeplitz-variation}
 \sum_{j=0}^{n}|w_{n,j}|\leq16.
\end{equation}
Every branch index sampled in the $n$th row tends uniformly to infinity, and
the real, unshifted spline branches converge to $\Th(x)$.  Therefore the
finite-row Toeplitz lemma, using
\eqref{eq:toeplitz-row-sum}--\eqref{eq:toeplitz-variation}, shows that their
weighted average tends to $\Th(x)$.

It remains to restore $q$.  On a fixed real branch, translation has the exact
finite Taylor expansion
\begin{equation}\label{eq:spline-translation}
 H_{p,M}(z+q)-H_{p,M}(z)
 =\sum_{d=0}^{p-1}
 \frac{2^{\binom d2}}{2^{\binom p2}(p-d)!}
 q^{p-d}H_{d,M}(z),
\end{equation}
where $H_{p,M}$ denotes the normalized degree-$p$ Thue--Morse branch with
prefix length $M$.  The lower real branches have norm at most one in the
relevant range, whence
\begin{equation}\label{eq:spline-translation-bound}
 |H_{p,M}(z+q)-H_{p,M}(z)|
 \leq2^{-(p-1)}\e^{|q|}.
\end{equation}
Since $p=2n-j\geq n$ in every row, this perturbation tends to zero uniformly
in $j$.  Combining it with the bounded total variation proves
\eqref{eq:discrete-limit} over both $\R$ and $\C$.
\end{proof}

\begin{remark}
The finite approximants $D_n(q,x)$ can depend on $q$; only their limit is
translation-independent.  This differs from
\Cref{thm:global-q-series}, where each fixed-scale polynomial is already
constant in $q$.
\end{remark}

\section{Concordance with the formal development and source corrections}
\label{app:concordance}

The article is organized mathematically rather than by Lean file.  The
following concordance indicates where the principal source groups have been
translated into ordinary analysis.

\begingroup
\small
\begin{longtable}{@{}>{\raggedright\arraybackslash}p{0.31\textwidth}>{\raggedright\arraybackslash}p{0.62\textwidth}@{}}
\toprule
Formal or printed source group & Human-readable treatment in this article\\
\midrule
\endfirsthead
\toprule
Formal or printed source group & Human-readable treatment in this article\\
\midrule
\endhead
Arias de Reyna, first paper: existence, uniqueness, convolution, and
probability model
& \Cref{sec:three-functions} and the construction in the sections beginning
with ``Construction by probability, convolution, and Fourier transform.''\\

First paper: polynomial step approximants and their pointwise limit
& \Cref{sec:basic-properties}, including the half-weight endpoint convention.\\

First paper: partition of unity, Poisson summation, cosine product, and
Weierstrass product
& The latter half of \Cref{sec:basic-properties}.\\

First paper: derivative identities, integer values, non-analyticity, and
moments
& \Cref{sec:global-extension} and the moment section beginning on
page~\pageref{sec:global-extension}.\\

First paper: rationality and the finite dyadic formula
& \Cref{thm:closed-dyadic} and its proof in the section ``A closed formula for
every dyadic value.''\\

Arithmetic paper: $c_n,d_n$, their generating functions and recurrences
& The moment section, especially \eqref{eq:c-recurrence},
\eqref{eq:f-generating}, and \eqref{eq:d-recurrence}.\\

Arithmetic paper: exact inverse powers, Taylor recursion, and the executable
evaluator
& \Cref{sec:bit-recursion}.\\

Arithmetic paper: integer normalizations, denominator bounds, parity,
valuations, and Bernoulli recurrences
& \Cref{sec:arithmetic}; the still open denominator equality is explicitly
kept as a conjecture.\\

Additional formal development: finite translated $q$-binomial formula
& \Cref{sec:qbinomial}.\\

Additional formal development: Fourier--Legendre coefficients and uniform
series
& \Cref{sec:legendre}.\\

Additional formal development: negative Laplace product, Mellin finite part,
Gamma--zeta fluctuation, exact Lambert phase, saddle method, and all orders
& The endpoint sequence of sections from ``The negative Laplace product''
through \Cref{sec:full-expansion}.\\

Additional formal development: binary telescope, corrected infinite
$q$-binomial series, parity-power series, and discrete limit
& \Cref{app:global-series}.\\
\bottomrule
\end{longtable}
\endgroup

\subsection{Corrections incorporated in the statements}

Several printed or informal formulas require small but mathematically
important repairs.  They have been incorporated silently into the main
statements and are listed here for auditability.

\begin{enumerate}[label=\arabic*.]
\item The convolution occurring in the first paper's equation~(12) is finite;
an infinite upper index is incompatible with the displayed dependence on the
stage parameter.

\item Closed interval indicators in the early step approximants double-count
shared endpoints.  The correct normalized approximants give each shared
endpoint weight $1/2$.

\item In the Poisson formulas, the exponential phase retains the translation
parameter, the differentiated identity requires positive order, and the
closed dyadic formula uses the consistent scaling adopted in
\Cref{thm:closed-dyadic}.

\item The bounded CDF-style function satisfies $F'(x)=2F(2x)$ only on its left
half.  The equation valid on all of $\R$ belongs to the signed extension
$\Th$.

\item The translated derivative lemma in the arithmetic paper requires the
nonnegative scale-plus-order hypothesis used in its proof.

\item In the integer $R_n$, the power of two is the positive exponent
consistent with the recurrence, proof, and displayed numerical values.

\item The naturality identity used in the parity argument contains the factor
of two forced by the normalization of the half moments.

\item The global binary and $q$-binomial series begin at $m=0$, with
$B_0=\lfloor x/2\rfloor$.  The omitted term vanishes on $[0,1)$ but equals one
at $x=1$.

\item A sharp endpoint formula cannot discard the Gamma--zeta periodic
fluctuation.  Its nonzero Fourier modes are present at the exact lower-Lambert
phase; replacing that phase prematurely also displaces terms involving
$\log\log(1/x)$.
\end{enumerate}

\begin{thebibliography}{99}

\bibitem{Arias1982}
J.~Arias de Reyna,
\emph{Definici\'on y estudio de una funci\'on indefinidamente diferenciable
de soporte compacto},
Rev. Real Acad. Cienc. Exact. F\'is. Natur. Madrid \textbf{76} (1982),
21--38; English translation,
\emph{An infinitely differentiable function with compact support: definition
and properties}, arXiv:1702.05442 (2017).

\bibitem{AriasArithmetic}
J.~Arias de Reyna,
\emph{Arithmetic of the Fabius function},
Integers \textbf{18} (2018), Paper A51; arXiv:1702.06487, version~3.

\bibitem{Corless}
R.~M. Corless, G.~H. Gonnet, D.~E.~G. Hare, D.~J. Jeffrey, and D.~E. Knuth,
\emph{On the Lambert $W$ function},
Adv. Comput. Math. \textbf{5} (1996), 329--359.

\bibitem{Fabius1966}
J.~Fabius,
\emph{A probabilistic example of a nowhere analytic $C^\infty$-function},
Z. Wahrscheinlichkeitstheorie verw. Geb. \textbf{5} (1966), 173--174.

\bibitem{Rvachev1990}
V.~A. Rvachev,
\emph{Compactly supported solutions of functional-differential equations and
their applications},
Russian Math. Surveys \textbf{45}:1 (1990), 87--120.

\bibitem{ProveIt}
V.~Reshetnikov,
\emph{ProveIt: FabiusFunction},
Lean formalization in the directory
\texttt{Analysis/FabiusFunction}, GitHub repository, accessed 24 August 2026.

\end{thebibliography}

\end{document}
